% MAT670 -- Interlude C: the torus, Diophantine approximation and aperiodic order.
% Entirely new material (not in the handwritten notes). Labels prefixed ERG:.
% Placed after Interlude B (number theory). Numerical claims checked with python3.

\tikzset{
  ERGv/.style={circle,draw,fill=white,inner sep=0pt,minimum size=6pt},
  ERGw/.style={circle,draw,fill=addcol!25,inner sep=0pt,minimum size=6pt},
  ERGproj/.style={gray,densely dotted},
  ERGthick/.style={draw=black,line width=0.3pt,fill=addcol!35},
  ERGthin/.style={draw=black,line width=0.3pt,fill=flagcol!30},
}

\chapter[Interlude C: Torus and aperiodic order]{Interlude C: The torus, Diophantine approximation and aperiodic order}
\label{ERG:ch}

\begin{addedblock}
This interlude is new: it is not in the handwritten notes. It starts from the torus $\R^n/\Lat$ of
Lecture~4 and asks what a single straight line does on it. The answer leads through Diophantine
approximation, the points on a circle of Lecture~7, the Fibonacci words of Lecture~11, and finally
to quasicrystals, which are obtained by \emph{folding} the Dynkin diagrams of Lectures~4--6. To keep
it readable it is set in black; all of it should be read as added material.
\end{addedblock}

Several objects in these notes have quietly been tori. Lecture~4 attached to a lattice $\Lat\subset\R^n$
the torus $\R^n/\Lat$, whose volume is $\abs{\det\Lat}$. Interlude~A found the maximal torus of
$\mathrm{SU}(n)$ to be $\mathfrak t/2\pi A_{n-1}$ (Section~\ref{LIE:sec:torus}). And the circle
$\Sph^1=\R/\Z$ carrying the points of Lecture~7 is the one-dimensional torus. The simplest dynamical
question we can ask about a torus is:
\begin{center}
\emph{draw a straight line on $\R^2/\Z^2$. Where does it go?}
\end{center}
If the slope is rational, the line closes up. If the slope is irrational, the line never closes and
\emph{shreds} the torus: it comes arbitrarily close to every point. How fast it does so is the subject
of Diophantine approximation, and the answer singles out the golden ratio
\[
  \varphi=\tau=\frac{1+\sqrt5}2=1.618\ldots,\qquad \varphi^2=\varphi+1 .
\]
(Both letters are standard. We write $\varphi$ in the number-theoretic sections and $\tau$ in the
Coxeter group sections, following each literature; $\tau'=-1/\tau=\frac{1-\sqrt5}{2}$ is its Galois
conjugate. Here $\tau$ is \emph{not} the kissing number.) Cutting the line out of the torus in a
different way produces the Fibonacci chain, the simplest quasicrystal. Its higher-dimensional
relatives come from the lattices $A_4$, $D_6$ and $E_8$ of Lecture~6.

Throughout, $\{x\}=x-\lfloor x\rfloor$ is the fractional part, $\norm{x}_{\Z}=\min_{p\in\Z}\abs{x-p}$
is the distance to the nearest integer, and $e(x)=e^{2\pi ix}$.

%======================================================================
\section{Shredding the torus}
\label{ERG:sec:torus}

Let $\ell_\alpha$ be the image in $\T^2=\R^2/\Z^2$ of the line $y=\alpha x$. Each time the line crosses
the vertical circle $\{x=0\}$ (that is, at $x=k\in\Z$) its height is $k\alpha$ modulo $1$. So the
\emph{first return map} to this circle is the rotation
\[
  R_\alpha\colon\T^1\to\T^1,\qquad t\mapsto t+\alpha \pmod 1,
\]
and $\ell_\alpha$ is the suspension of $R_\alpha$: it is dense in $\T^2$ exactly when the orbit
$\{k\alpha\}$ is dense in $\T^1$.

\begin{figure}[ht]
\centering
\begin{tikzpicture}[scale=2.5]
  \begin{scope}
    \begin{scope}\clip (0,0) rectangle (1,1);
      \foreach \k in {0,...,4} {\pgfmathsetmacro{\y}{\k*0.4-floor(\k*0.4+0.0001)}
        \draw[addcol,thick] (0,\y)--++(1,0.4); \draw[addcol,thick] (0,\y-1)--++(1,0.4);}
    \end{scope}
    \draw[thick] (0,0) rectangle (1,1);
    \foreach \k in {0,...,4} {\pgfmathsetmacro{\y}{\k*0.4-floor(\k*0.4+0.0001)} \fill (0,\y) circle (0.025);}
    \node[below,font=\footnotesize] at (0.5,-0.05) {slope $\frac25$: a closed curve};
  \end{scope}
  \begin{scope}[shift={(1.8,0)}]
    \begin{scope}\clip (0,0) rectangle (1,1);
      \foreach \k in {0,...,29} {\pgfmathsetmacro{\y}{\k*0.41421356-floor(\k*0.41421356)}
        \draw[addcol] (0,\y)--++(1,0.41421356); \draw[addcol] (0,\y-1)--++(1,0.41421356);}
    \end{scope}
    \draw[thick] (0,0) rectangle (1,1);
    \foreach \k in {0,...,29} {\pgfmathsetmacro{\y}{\k*0.41421356-floor(\k*0.41421356)} \fill (0,\y) circle (0.018);}
    \node[below,font=\footnotesize] at (0.5,-0.05) {slope $\sqrt2-1$: $30$ windings};
  \end{scope}
  \foreach \s in {0,1.8} {
    \draw[-{Stealth[length=3pt]}] (\s+0.47,0)--(\s+0.53,0); \draw[-{Stealth[length=3pt]}] (\s+0.47,1)--(\s+0.53,1);
    \draw[-{Stealth[length=3pt]}] (\s,0.44)--(\s,0.5); \draw[-{Stealth[length=3pt]}] (\s,0.5)--(\s,0.56);
    \draw[-{Stealth[length=3pt]}] (\s+1,0.44)--(\s+1,0.5); \draw[-{Stealth[length=3pt]}] (\s+1,0.5)--(\s+1,0.56);}
\end{tikzpicture}
\caption{The line $y=\alpha x$ wrapped around the torus $\R^2/\Z^2$ (opposite edges identified as
marked). The dots on the left edge are the returns $\{k\alpha\}$ to the circle $x=0$. For rational
$\alpha=p/q$ the line closes after $q$ turns; for irrational $\alpha$ it shreds the torus.}
\label{ERG:fig:torus}
\end{figure}

If $\alpha=p/q$ in lowest terms, the line closes up after $q$ horizontal and $p$ vertical turns, a
closed geodesic of length $\sqrt{p^2+q^2}$. The irrational case is Kronecker's theorem, and the proof
is the pigeonhole principle.

\begin{proposition}[Dirichlet 1842; Kronecker 1884]\label{ERG:prop:dirichlet}
Let $\alpha\in\R$ and let $Q\ge1$ be an integer. There is an integer $q$ with $1\le q\le Q$ and
$\norm{q\alpha}_\Z<1/Q$. If $\alpha$ is irrational, the orbit $\{k\alpha\}_{k\ge0}$ is dense in
$[0,1)$, and so $\ell_\alpha$ is dense in $\T^2$.
\end{proposition}

\begin{proof}
The $Q+1$ numbers $\{k\alpha\}$, $0\le k\le Q$, lie in the $Q$ intervals $[i/Q,(i+1)/Q)$, so two of
them, say for $j<k$, lie in the same interval. Then $q=k-j$ satisfies $\norm{q\alpha}_\Z<1/Q$. If
$\alpha$ is irrational then $\delta=\norm{q\alpha}_\Z>0$, and the multiples of $q\alpha$ march around
the circle in steps of length $\delta<1/Q$ (all in the same direction), so every interval of length
$1/Q$ contains a point of the orbit. Since $Q$ was arbitrary, the orbit is dense.
\end{proof}

Density says nothing about \emph{how often} the orbit visits a given arc. The right notion is the one
behind the Boltzmann distribution of Lecture~7: time averages should equal space averages.

\begin{definition}
A sequence $(x_k)$ in $[0,1)$ is \emph{equidistributed} if for every interval $I\subseteq[0,1)$
\[
  \lim_{N\to\infty}\frac1N\#\{0\le k<N: x_k\in I\}=\abs{I}.
\]
\end{definition}

\begin{theorem}[Weyl's criterion; Weyl 1916]\label{ERG:thm:weylcrit}
$(x_k)$ is equidistributed if and only if
\[
  \frac1N\sum_{k<N}e(mx_k)\to0\qquad\text{for every integer } m\neq0.
\]
Equivalently, $\frac1N\sum_{k<N}f(x_k)\to\int_0^1f$ for every continuous $1$-periodic $f$.
\end{theorem}

\begin{proof}[Sketch]
If $(x_k)$ is equidistributed, then the averages of step functions converge, and continuous functions are
uniform limits of step functions. Conversely, the condition gives convergence for trigonometric
polynomials, which are uniformly dense in the continuous periodic functions (Fej\'er, or
Stone--Weierstrass). For an interval $I$, choose continuous $f_-\le\mathbf 1_I\le f_+$ with
$\int(f_+-f_-)<\varepsilon$ and squeeze.
\end{proof}

\begin{theorem}[Weyl 1916; also Bohl and Sierpi\'nski 1909--10]\label{ERG:thm:weyl}
If $\alpha$ is irrational, $(\{k\alpha\})_{k\ge0}$ is equidistributed.
\end{theorem}

\begin{proof}
For $m\neq0$ we have $m\alpha\notin\Z$, so $e(m\alpha)\neq1$ and the geometric series gives
\[
  \Big|\frac1N\sum_{k=0}^{N-1}e(mk\alpha)\Big|
  =\frac{\abs{e(mN\alpha)-1}}{N\abs{e(m\alpha)-1}}\le\frac{1}{N\abs{\sin\pi m\alpha}}\longrightarrow0 .
  \qedhere
\]
\end{proof}

The same Fourier argument gives a statement about measures rather than orbits.

\begin{proposition}[Unique ergodicity]\label{ERG:prop:ue}
If $\alpha$ is irrational, Lebesgue measure is the only Borel probability measure on $\T^1$ invariant
under $R_\alpha$.
\end{proposition}

\begin{proof}
Let $\mu$ be invariant and $\hat\mu(m)=\int e(-mt)\,d\mu(t)$. Invariance gives
$\hat\mu(m)=\int e(-m(t+\alpha))\,d\mu=e(-m\alpha)\hat\mu(m)$, so $\hat\mu(m)=0$ for $m\neq0$. A
measure on the circle is determined by its Fourier coefficients, and these are those of Lebesgue measure.
\end{proof}

For a general measure-preserving system there is no such uniformity, and the best one can say is
\emph{Birkhoff's ergodic theorem} (1931): if $T$ preserves a probability measure $\mu$ and
$f\in L^1(\mu)$, then $\frac1N\sum_{k<N}f(T^kx)$ converges for $\mu$-almost every $x$, and the limit is
$\int f\,d\mu$ when $T$ is \emph{ergodic} (every invariant set has measure $0$ or $1$). Unique
ergodicity upgrades ``almost every $x$'' to ``every $x$, uniformly''.

\subsection*{Higher dimensions, and the torus inside a Lie group}

Now let $\Lat\subset\R^n$ be a lattice and $\T_\Lat=\R^n/\Lat$ the torus of Lecture~4. Its
characters are $x\mapsto e(\ip{w}{x})$ with $w$ in the \emph{dual lattice}
$\Lat^*=\{w:\ip{w}{\lambda}\in\Z\ \forall\lambda\in\Lat\}$ (the same dual lattice that appears in the
Poisson summation formula of Interlude~B). Weyl's criterion holds on $\T_\Lat$ with these characters, and
the geometric series argument gives:

\begin{theorem}[Kronecker--Weyl]\label{ERG:thm:kroneckerweyl}
Let $v\in\R^n$.
\begin{enumerate}[nosep]
\item The sequence $kv\bmod\Lat$ ($k\in\N$) is equidistributed in $\T_\Lat$ if and only if
  $\ip{w}{v}\notin\Z$ for every $w\in\Lat^*\setminus\{0\}$.
\item The line $tv\bmod\Lat$ ($t\in\R$) is equidistributed (in particular dense) if and only if
  $\ip{w}{v}\neq0$ for every $w\in\Lat^*\setminus\{0\}$, i.e.\ $v$ lies in no rational hyperplane of
  $\Lat$.
\end{enumerate}
For $\Lat=\Z^n$, (1) says that $1,v_1,\dots,v_n$ are linearly independent over $\Q$.
\end{theorem}

If the condition fails for some $w$, the orbit stays in the closed subgroup $\{x:\ip{w}{x}\in\Z\}$,
a finite union of parallel subtori. So orbit closures
of linear flows on tori are always subtori. This is the abelian ancestor of Ratner's theorem in
Section~\ref{ERG:sec:lattices}.

\begin{remark}[Generic elements of a maximal torus]
In Interlude~A the maximal torus of a compact group is $T=\mathfrak t/\Lat_0$, with $\Lat_0$ the kernel
of $\exp$ (Section~\ref{LIE:sec:torus}). By Theorem~\ref{ERG:thm:kroneckerweyl}(1), an element
$t=\exp(H)$ generates a dense subgroup of $T$ if and only if $\chi(t)\neq1$ for every nontrivial
character $\chi$ of $T$. Such \emph{topological generators} form a dense set of full measure. They are
the standard tool for deducing from Theorem~\ref{LIE:thm:torus} (every element is conjugate into $T$)
that any two maximal tori are conjugate: conjugate a generator $t'$ of $T'$ into $T$, and its closure
$T'$ goes along with it.
\end{remark}

%======================================================================
\section{How close does the line come? Diophantine approximation}
\label{ERG:sec:dioph}

The line $\ell_\alpha$ passes within $\norm{q\alpha}_\Z$ of the lattice point $(q,p)$, measured
vertically. Dirichlet's bound $\norm{q\alpha}_\Z<1/q$ for infinitely many $q$ is a statement about
lattice points near a line, and it has a second proof from the geometry of numbers of Lecture~4.

\begin{proof}[Second proof of Dirichlet, via Minkowski]
The parallelogram $K_Q=\{(x,y):\abs{x}\le Q,\ \abs{x\alpha-y}\le1/Q\}$ is convex, symmetric and compact,
of area $2Q\cdot 2/Q=4$. By Minkowski's convex body theorem it contains a lattice point
$(q,p)\neq0$. If $q=0$ then $\abs{p}\le1/Q<1$ forces $p=0$; so $1\le\abs{q}\le Q$ and
$\abs{q\alpha-p}\le1/Q$.
\end{proof}

\subsection*{Continued fractions}

Write $\alpha=[a_0;a_1,a_2,\dots]=a_0+\cfrac{1}{a_1+\cfrac{1}{a_2+\dotsb}}$ with $a_i\in\Z$, $a_i\ge1$
for $i\ge1$, and let $p_n/q_n=[a_0;a_1,\dots,a_n]$ be the convergents, so
$q_{n+1}=a_{n+1}q_n+q_{n-1}$. The following facts are classical (see e.g.\ Khinchin's
\emph{Continued Fractions}, or Cassels' \emph{Introduction to Diophantine Approximation}, Ch.~I).
\begin{enumerate}[nosep,label=(\roman*)]
\item $\frac{1}{(a_{n+1}+2)}<q_n\,\norm{q_n\alpha}_\Z<\frac{1}{a_{n+1}}$ for $n\ge1$, and $q_n\alpha-p_n$
  alternates in sign.
\item (Best approximation, Lagrange) If $1\le q<q_{n+1}$, then $\norm{q\alpha}_\Z\ge\norm{q_n\alpha}_\Z$.
  So the record holders for $\norm{q\alpha}_\Z$ are exactly the $q_n$.
\item For $\varphi=[1;1,1,1,\dots]$ the convergents are $F_{n+1}/F_n$ (Fibonacci numbers,
  $F_1=F_2=1$), and $F_n\varphi-F_{n+1}=-(\tau')^{n}=(-1)^{n+1}\varphi^{-n}$.
\end{enumerate}
Fact (iii) is Binet's formula $\varphi^n=F_n\varphi+F_{n-1}$ applied to $\tau'$, which satisfies the
same quadratic equation. By (i), $\inf_q q\norm{q\alpha}_\Z>0$ if and only if the partial quotients $a_n$
are bounded; such $\alpha$ are called \emph{badly approximable}. The golden ratio has the smallest
possible partial quotients, and it turns out to be the worst approximable number of all.

\begin{theorem}[Hurwitz 1891]\label{ERG:thm:hurwitz}
For every irrational $\alpha$ there are infinitely many fractions $p/q$ with
$\abs{\alpha-p/q}<\frac1{\sqrt5\,q^2}$. The constant $\sqrt5$ cannot be increased: for
$\alpha=\varphi$ and any $A>\sqrt5$ only finitely many $p/q$ satisfy $\abs{\alpha-p/q}<\frac1{Aq^2}$.
\end{theorem}

\begin{proof}[Proof of the sharpness]
The existence part follows by showing that of any three consecutive convergents one satisfies the
inequality (Hardy--Wright, Thm.~195). For the optimality, note that
$(p-q\varphi)(p-q\tau')=p^2-pq-q^2$ is a nonzero integer for $(p,q)\neq(0,0)$, because $\varphi$ is
irrational. Hence
\[
  q^2\Big|\varphi-\frac pq\Big|=\frac{q\abs{p^2-pq-q^2}}{\abs{p-q\tau'}}\ \ge\ \frac{1}{\abs{p/q-\tau'}}
  \longrightarrow\frac1{\varphi-\tau'}=\frac1{\sqrt5}
\]
as $p/q\to\varphi$. So $\liminf q^2\abs{\varphi-p/q}\ge1/\sqrt5$, and the convergents
$F_{n+1}/F_n$ attain this value in the limit.
\end{proof}

The proof is really about the indefinite quadratic form $x^2-xy-y^2$, which takes no small nonzero
values on $\Z^2$. Markov turned this into a complete description of the top of the spectrum.

\begin{theorem}[Markov 1879--80]\label{ERG:thm:markov}
For a real indefinite binary quadratic form $f$ of discriminant $d>0$, put
$\mu(f)=\inf_{(x,y)\in\Z^2\setminus0}\abs{f(x,y)}/\sqrt d$. The values of $\mu(f)$ that exceed $\frac13$
are exactly
\[
  \frac{m}{\sqrt{9m^2-4}},\qquad m\in\{1,2,5,13,29,34,89,169,194,\dots\},
\]
where $m$ runs over the \emph{Markov numbers}, the integers occurring in positive solutions of
$x^2+y^2+z^2=3xyz$. The value for $m$ is attained only by forms equivalent to a multiple of an
explicit integral form $f_m$. For $m=1,2$ these are $x^2-xy-y^2$ ($\mu=1/\sqrt5$) and $x^2-2y^2$
($\mu=1/\sqrt8$).
\end{theorem}

The dual statement for numbers is that the \emph{Lagrange spectrum}
$\{\limsup_q(q\norm{q\alpha}_\Z)^{-1}\}$ begins $\sqrt5,\sqrt8,\sqrt{221}/5,\dots$ and accumulates
only at $3$. All Markov triples arise from $(1,1,1)$ by the moves $(x,y,z)\mapsto(x,y,3xy-z)$. Whether
each Markov number is the largest entry of only one triple (Frobenius 1913) is the \emph{Markov
uniqueness conjecture}, still open.

\subsection*{Badly approximable numbers are bounded orbits}

Here the torus meets the space of lattices $X_2=\SL_2(\R)/\SL_2(\Z)$ of Lecture~4. Put
\[
  \Lat_\alpha=\begin{pmatrix}1&0\\ \alpha&1\end{pmatrix}\Z^2=\{(q,\,q\alpha+p):p,q\in\Z\},
  \qquad g_t=\begin{pmatrix}e^{-t}&0\\0&e^{t}\end{pmatrix}.
\]
The flow $g_t$ squeezes the first coordinate and stretches the second. A lattice point with small
$q\abs{q\alpha+p}$ becomes, at the right time $t$, a short vector of $g_t\Lat_\alpha$.

\begin{proposition}[Dani correspondence, simplest case]\label{ERG:prop:dani}
$\alpha$ is badly approximable if and only if the forward orbit $\{g_t\Lat_\alpha:t\ge0\}$ is bounded
in $X_2$.
\end{proposition}

\begin{proof}
By Mahler's criterion (Lecture~4) a set of unimodular lattices is bounded exactly when the shortest
vector length $\lambda_1$ is bounded below on it. We use the sup norm, which changes $\lambda_1$ by at
most a factor $\sqrt2$. Suppose $\abs q\abs{q\alpha+p}\ge c$ for all $q\neq0$, where we may take $c\le1$.
A vector $v=(e^{-t}q,e^t(q\alpha+p))$ of $g_t\Lat_\alpha$ with $q=0$ has $\norm v_\infty=e^t\abs p\ge1$.
If $q\neq0$ then $\norm v_\infty^2\ge\abs q\abs{q\alpha+p}\ge c$. So $\lambda_1\ge\sqrt c$ along the
orbit. Conversely, suppose $\norm v_\infty\ge\varepsilon$ for all $t\ge0$ and $v\neq0$, where
$\varepsilon\le1$. Given $q\ge1$ and $p$, choose $t\ge0$ with $e^{-t}q=\varepsilon/2$. Then the first
coordinate of $v$ is less than $\varepsilon$, so $e^t\abs{q\alpha+p}\ge\varepsilon$, that is,
$q\abs{q\alpha+p}\ge\varepsilon^2/2$.
\end{proof}

This dictionary, due to Dani (1985) in all dimensions, turns questions about numbers into questions about
orbits on the space of lattices. We return to it in Section~\ref{ERG:sec:lattices}.

%======================================================================
\section{Three gaps}
\label{ERG:sec:threegap}

Return to the circle of Lecture~7 (Section~\ref{L07:subsec:circle}), which we take to have length $1$, and
place on it the first $N$ points of an irrational rotation,
\[
  x_k=\{k\alpha\},\qquad 0\le k\le N-1 .
\]
They cut the circle into $N$ arcs, a point $(g_1,\dots,g_N)$ of the open gap simplex of Lecture~7. One
would expect $N$ different gap lengths. There are at most three.

\begin{theorem}[Three gap theorem; conjectured by Steinhaus, proved by S\'os 1957/58, Sur\'anyi 1958, \'Swierczkowski 1959]\label{ERG:thm:threegap}
Let $\alpha$ be irrational and $N\ge2$. Let $u,v\in\{1,\dots,N-1\}$ be the indices of the points
closest to $0$ from the right and from the left, and put $a=\{u\alpha\}$, $b=1-\{v\alpha\}$. Then every
gap between neighbouring points has length $a$, $b$ or $a+b$. More precisely, the right neighbour of
$x_k$ is
\[
  x_{k+u}\ \ (\text{gap } a)\ \text{if } k<N-u,\qquad
  x_{k-v}\ \ (\text{gap } b)\ \text{if } k\ge v,\qquad
  x_{k+u-v}\ \ (\text{gap } a+b)\ \text{otherwise}.
\]
There are $N-u$ gaps of length $a$, $N-v$ of length $b$, and $u+v-N$ of length $a+b$.
\end{theorem}

\begin{proof}
All the indices named are in $[0,N-1]$, so each claimed neighbour is one of the points, at the stated
distance to the right of $x_k$ (because $x_{k+u}=x_k+a$ and $x_{k-v}=x_k+b$ modulo $1$). We must
show that no point $x_j$ lies strictly inside the arc. We use only the minimality of $a$ and $b$: no
$1\le m\le N-1$ has $\{m\alpha\}\in(0,a)$ or $\{m\alpha\}\in(1-b,1)$.

\emph{Gap $a$} ($k+u\le N-1$). If $x_j\in(x_k,x_k+a)$ with $j>k$, then $\{(j-k)\alpha\}\in(0,a)$,
contradicting minimality of $a$. If $j<k$, then $x_{k+u}=x_k+a$ lies in $(x_j,x_j+a)$, so
$m=k+u-j\in[1,N-1]$ has $\{m\alpha\}\in(0,a)$, again a contradiction.

\emph{Gap $b$} ($k\ge v$). If $x_j\in(x_k,x_k+b)$ with $j<k$, then $\{(j-k)\alpha\}\in(0,b)$, i.e.\
$m=k-j$ has $\{m\alpha\}\in(1-b,1)$. If $j>k$, then $x_j\in(x_{k-v}-b,x_{k-v})$, and $m=j-k+v$
(which lies in $[1,N-1]$) has $\{m\alpha\}\in(1-b,1)$. Both contradict minimality of $b$.

\emph{Gap $a+b$} ($N-u\le k\le v-1$). Suppose $x_j\in(x_k,x_k+a+b)$. If $j>k$, then
$\{(j-k)\alpha\}\in(0,a+b)$, and in fact in $(a,a+b)$ by minimality of $a$. (It cannot equal $a$,
since then $j=k+u\ge N$.) Then $m=k+u-j$ lies in $[1,N-1]$ and $\{-m\alpha\}=\{(j-k)\alpha\}-a\in(0,b)$, so
$\{m\alpha\}\in(1-b,1)$, which is impossible. If $j<k$, then $s=k-j\in[1,N-1]$ has
$1-\{s\alpha\}\in(b,a+b)$, and $m=v-s\in[1,N-1]$ has
$\{m\alpha\}=\{v\alpha\}-\{s\alpha\}=(1-\{s\alpha\})-b\in(0,a)$, which is impossible.

The three ranges of $k$ are disjoint: a $k$ in the first two would give $a=b$, hence
$(u+v)\alpha\in\Z$. Counting the $k$ in each range gives the multiplicities.
\end{proof}

\begin{figure}[ht]
\centering
\begin{tikzpicture}[scale=1.75]
  \draw[gray!40] (0,0) circle (1);
  % arcs: start x, gap, style
  \foreach \x/\g/\c in {0/0.0711/addcol, 0.0711/0.1716/black, 0.2426/0.0711/addcol, 0.3137/0.1005/flagcol,
     0.4142/0.0711/addcol, 0.4853/0.1716/black, 0.6569/0.0711/addcol, 0.7279/0.1005/flagcol,
     0.8284/0.0711/addcol, 0.8995/0.1005/flagcol}
    {\draw[\c,line width=1.8pt] ({90-360*\x}:1) arc[start angle={90-360*\x}, end angle={90-360*(\x+\g)}, radius=1];}
  \foreach \k/\x in {0/0, 5/0.0711, 3/0.2426, 8/0.3137, 1/0.4142, 6/0.4853, 4/0.6569, 9/0.7279, 2/0.8284, 7/0.8995}
    {\fill ({90-360*\x}:1) circle (0.035); \node[font=\scriptsize] at ({90-360*\x}:1.17) {$\k$};}
  \node[font=\small,align=left,anchor=west] at (1.45,0.35) {\textcolor{addcol}{$\blacksquare$} $a=\{5\alpha\}\approx0.0711$ \ ($\times5$)};
  \node[font=\small,align=left,anchor=west] at (1.45,0) {\textcolor{flagcol}{$\blacksquare$} $b=1-\{7\alpha\}\approx0.1005$ \ ($\times3$)};
  \node[font=\small,align=left,anchor=west] at (1.45,-0.35) {$\blacksquare$ $a+b\approx0.1716$ \ ($\times2$)};
\end{tikzpicture}
\caption{The points $\{k\alpha\}$, $0\le k\le9$, for $\alpha=\sqrt2-1$, drawn clockwise from the top.
Here $u=5$, $v=7$, $N=10$, and the multiplicities are $N-u=5$, $N-v=3$ and $u+v-N=2$.}
\label{ERG:fig:threegap}
\end{figure}

\begin{example}[Fibonacci: two gaps in ratio $\varphi$]\label{ERG:ex:fibgaps}
Take $\alpha=\varphi^{-1}=\varphi-1$ and $N=F_n$ with $n\ge4$. By facts (ii) and (iii) above, the
closest returns to $0$ among $1\le k<F_n$ occur at $k=F_{n-1}$ and $k=F_{n-2}$. (Since all partial
quotients equal $1$, there are no intermediate fractions, so the one-sided best approximations are also
convergents.) They are at distances
$\varphi^{-(n-1)}$ and $\varphi^{-(n-2)}$ on opposite sides (the convergents of $\varphi^{-1}$ are
$F_{k-1}/F_k$). So $\{u,v\}=\{F_{n-1},F_{n-2}\}$ and $u+v-N=0$: \emph{there are only two gap lengths,
and their ratio is $\varphi$}. There are $F_{n-2}$ short gaps and $F_{n-1}$ long ones, and indeed
$F_{n-2}\varphi^{-(n-1)}+F_{n-1}\varphi^{-(n-2)}=1$. Adding the next point splits one long gap into a
long and a short gap, in the ratio $\varphi:1$. This is the arrangement of leaves (``phyllotaxis'') that
a plant produces when it turns each new leaf by the golden angle $2\pi/\varphi^2\approx137.5^\circ$.
(Checked numerically for $N\le40$.)
\end{example}

\begin{proposition}[How well spread is an orbit?]\label{ERG:prop:spread}
Let $m_N(\alpha)$ be the smallest gap of $\{k\alpha\}_{0\le k<N}$, so that the injectivity radius of
Lecture~7 is $\rho=m_N/2$. The regular $N$-gon has $Nm_N=1$. For every irrational $\alpha$,
\[
  \liminf_{N\to\infty}N\,m_N(\alpha)=\liminf_{q\to\infty}q\,\norm{q\alpha}_\Z\le\frac1{\sqrt5},
\]
with equality for $\alpha=\varphi$. Moreover $\inf_NNm_N(\alpha)>0$ if and only if $\alpha$ is badly
approximable.
\end{proposition}

\begin{proof}
The smallest gap is the smallest circular distance between two of the points, so
$m_N=\min_{1\le k<N}\norm{k\alpha}_\Z$. Taking $N=q+1$ gives
$Nm_N\le(q+1)\norm{q\alpha}_\Z$. Conversely $m_N=\norm{k\alpha}_\Z$ for some $k<N$, so
$Nm_N\ge k\norm{k\alpha}_\Z$, and $k\to\infty$ as $N\to\infty$ because $m_N\to0$. This proves the
equality of the two $\liminf$s, and the same two inequalities give the last statement. The bound
$1/\sqrt5$ and the equality case are Hurwitz's Theorem~\ref{ERG:thm:hurwitz}.
\end{proof}

So for \emph{every} rotation, infinitely often the injectivity radius falls short of the optimal value
$1/(2N)$ of Lecture~7 by a factor of (almost) $\sqrt5$. The golden rotation does no worse than this
asymptotically. Rotations by numbers not equivalent to $\varphi$ do worse (by a factor $\sqrt8$ or
more), by Markov's Theorem~\ref{ERG:thm:markov}.

\begin{remark}[Gaps as a function on the space of lattices]
The vector of gaps lives in the gap simplex of Lecture~7. In Interlude~A the same simplex appeared as
the Weyl alcove $T/W$ of $\mathrm{SU}(N)$ (the eigenvalues of a unitary matrix are $N$ points on a
circle). Marklof and Str\"ombergsson (2017) observed that the gap lengths of $\{k\alpha\}_{k<N}$ are
values of a simple function on the space of lattices $X_2$, evaluated at a lattice obtained from
$\Lat_\alpha$ by the diagonal flow at time $t\approx\log N$. The three gap theorem then becomes an elementary
statement about two-dimensional lattices.
\end{remark}

%======================================================================
\section{Fibonacci points and discrepancy}
\label{ERG:sec:fib}

How uniform can $N$ points in the unit square be? Measure the error with the \emph{star discrepancy}
\[
  D^*(P)=\sup_{0\le x,y\le1}\Big|\frac{\#\big(P\cap[0,x)\times[0,y)\big)}{N}-xy\Big| .
\]
Random points have $D^*\approx N^{-1/2}$. The three gap structure suggests a better choice. The
\emph{Fibonacci lattice} with $N=F_m$ points is
\[
  P_m=\Big\{\Big(\frac{k}{F_m},\Big\{\frac{kF_{m-1}}{F_m}\Big\}\Big):0\le k<F_m\Big\},
\]
the points of the line of slope $F_{m-1}/F_m\approx\varphi^{-1}$ on the torus
(Figure~\ref{ERG:fig:torus}), sampled at the $F_m$ times it crosses $x=k/F_m$. It is a lattice in
$\T^2$, namely $\Z^2+\Z\,(1,F_{m-1})/F_m$.

\begin{figure}[ht]
\centering
\begin{tikzpicture}[scale=3.4]
  \draw (0,0) rectangle (1,1);
  \foreach \k in {0,...,54} {\pgfmathsetmacro{\y}{mod(34*\k,55)/55} \fill[addcol] ({\k/55},\y) circle (0.012);}
  \node[below,font=\small] at (0.5,0) {$F_{10}=55$ Fibonacci lattice points};
\end{tikzpicture}
\qquad
\begin{tikzpicture}[scale=1.55]
  \draw[gray!50] (0,0) circle (1);
  \foreach \k in {0,...,99} {
    \pgfmathsetmacro{\z}{1-(2*\k+1)/100}
    \pgfmathsetmacro{\r}{sqrt(1-\z*\z)}
    \pgfmathsetmacro{\th}{360*mod(\k*0.6180339887,1)}
    \pgfmathsetmacro{\c}{cos(\th)}
    \ifdim \c pt>0pt \fill[addcol] ({\r*sin(\th)},\z) circle (0.022); \fi}
  \node[below,font=\small] at (0,-1.05) {$100$ spherical Fibonacci points};
\end{tikzpicture}
\caption{Left: the Fibonacci lattice $P_{10}$. Right: the front half of the spherical Fibonacci point set
\eqref{ERG:eq:sphfib} with $N=100$, viewed from the side.}
\label{ERG:fig:fib}
\end{figure}

\begin{theorem}[Discrepancy in the square]\label{ERG:thm:discrepancy}\leavevmode
\begin{enumerate}[nosep]
\item (Schmidt 1972) There is $c>0$ such that every $N$-point set $P\subset[0,1]^2$ has
  $D^*(P)\ge c\,\frac{\log N}{N}$. (Roth 1954 had proved the weaker bound $c\sqrt{\log N}/N$. It is a bound for the
  $L^2$ average of the discrepancy, and in that form it is sharp, by Davenport 1956.)
\item The Fibonacci lattice has $D^*(P_m)=O\big(\frac{\log F_m}{F_m}\big)$. More generally, a
  two-dimensional lattice $\{(k/N,\{kg/N\})\}$ has $D^*=O(K\log N/N)$ if the partial quotients of
  $g/N$ are at most $K$ (see Niederreiter, \emph{Random Number Generation and Quasi-Monte Carlo
  Methods}, 1992, Ch.~5).
\end{enumerate}
\end{theorem}

So the Fibonacci lattice is optimal up to a constant, because all partial quotients of
$F_{m-1}/F_m=[0;1,1,\dots,1]$ equal $1$. The same mechanism appears in one dimension: for badly
approximable $\alpha$ the discrepancy of $\{k\alpha\}_{k<N}$ is $O(\log N/N)$ (see Kuipers--Niederreiter,
Ch.~2), while Schmidt showed that every infinite sequence in $[0,1)$ has discrepancy at least $c\log N/N$ for infinitely many $N$.

\paragraph{On the sphere.} Pull the Fibonacci lattice back to $\Sph^2$ by the area-preserving
cylindrical (Lambert) projection $(\theta,z)\mapsto(\sqrt{1-z^2}\cos\theta,\sqrt{1-z^2}\sin\theta,z)$.
Up to small variations this gives the \emph{spherical Fibonacci points}
\begin{equation}\label{ERG:eq:sphfib}
  z_k=1-\frac{2k+1}{N},\qquad \theta_k=2\pi k\varphi^{-1},\qquad 0\le k<N,
\end{equation}
a golden spiral from pole to pole. The relevant measure of uniformity is the \emph{spherical cap
discrepancy} $D_{\mathrm{cap}}$, the supremum over spherical caps $C$ of $\abs{\#(P\cap C)/N-\sigma(C)}$,
with $\sigma$ the normalized area. Beck (1984) showed that every $N$-point set has
$D_{\mathrm{cap}}\ge cN^{-3/4}$, and that some $N$-point sets have $D_{\mathrm{cap}}\le C N^{-3/4}\sqrt{\log N}$.
For the spherical Fibonacci lattice, Aistleitner, Brauchart and Dick (2012) proved
$D_{\mathrm{cap}}\le 44\sqrt{8/F_m}$, that is $O(N^{-1/2})$. Their numerical experiments suggest
that the true order is close to $N^{-3/4}$, but this is \emph{not} proved. Spherical Fibonacci points
are also a standard starting configuration for numerical searches in the Tammes and Thomson problems
of Lectures~2 and~11. This is an empirical observation; they are not optimal.

%======================================================================
\section{Sturmian words and the Fibonacci chain}
\label{ERG:sec:sturm}

Now \emph{code} the irrational line. Record, at each integer time $k$, whether the rotation
$t\mapsto t+\alpha$ wraps around:
\[
  s_k=\lfloor(k+1)\alpha+\rho\rfloor-\lfloor k\alpha+\rho\rfloor\in\{0,1\}\qquad(0<\alpha<1).
\]
Equivalently, $s_k=1$ if and only if $\{k\alpha+\rho\}\in[1-\alpha,1)$. Such a word is called a
\emph{mechanical} (or \emph{Sturmian}) word of slope $\alpha$. It is also the \emph{cutting sequence} of
the line $y=\alpha x+\rho$: the record of which horizontal grid lines it crosses between consecutive
vertical ones. The \emph{complexity} $p(n)$ of an infinite word is its number of distinct factors (blocks)
of length $n$.

\begin{theorem}[Morse--Hedlund 1938/1940; Coven--Hedlund 1973]\label{ERG:thm:morsehedlund}
An infinite word with $p(n)\le n$ for some $n$ is eventually periodic. The one-sided infinite words of
the smallest remaining complexity, $p(n)=n+1$ for all $n$, are exactly the mechanical words of
irrational slope. (For two-sided words this needs recurrence: $\cdots aaabaaa\cdots$ also has
$p(n)=n+1$.)
\end{theorem}

\begin{proof}[Sketch that a mechanical word has $p(n)=n+1$]
The factor $s_k\cdots s_{k+n-1}$ is determined by which of the arcs $[1-\alpha,1)-j\alpha$,
$0\le j<n$, contain $y=\{k\alpha+\rho\}$. As $y$ moves around the circle, the factor can change only when
$y$ crosses one of the endpoints of these arcs, which are the $n+1$ points $\{-j\alpha\}$, $0\le j\le n$.
These cut the circle into $n+1$ arcs, so $p(n)\le n+1$. The word is not eventually periodic (its frequency
of $1$'s is the irrational number $\alpha$), so $p(n)\ge n+1$ by the first half of the theorem
(Exercise~\ref{ERG:ex:mh}). Hence $p(n)=n+1$, and the $n+1$ arcs correspond bijectively to the factors.
\end{proof}

By equidistribution, the frequency of a factor is the length of its arc, and by the three gap theorem
the $n+1$ arcs have at most three different lengths. So \emph{the factors of length $n$ of a Sturmian
word occur with at most three different frequencies}.

\paragraph{The Fibonacci word.} Let $\sigma$ be the substitution $a\mapsto ab$, $b\mapsto a$. Then
$\sigma^n(a)$ is a prefix of $\sigma^{n+1}(a)$, and the limit is the \emph{Fibonacci word}
\[
  f=abaababaabaababaababaabaababaabaab\cdots,\qquad \abs{\sigma^n(a)}=F_{n+2}.
\]
The lengths obey $\abs{\sigma^{n+1}(a)}=\abs{\sigma^n(a)}+\abs{\sigma^{n-1}(a)}$, because
$\sigma^{n+1}(a)=\sigma^n(a)\,\sigma^{n-1}(a)$. The letter $b$ never appears twice in a row, since
every $b$ comes from an $a\mapsto ab$ and the next block begins with $a$. The frequency of $b$ is
$\lim F_n/F_{n+2}=\varphi^{-2}\approx0.382$. In fact $f$ is the mechanical word of slope
$\varphi^{-2}$ with $\rho=\varphi^{-2}$ and $b\leftrightarrow1$ (the word with $\rho=0$ is $af$): its $k$-th letter ($k\ge0$) is $b$ exactly when
$\lfloor(k+2)\varphi^{-2}\rfloor-\lfloor(k+1)\varphi^{-2}\rfloor=1$ (Lothaire,
\emph{Algebraic Combinatorics on Words}, Ch.~2; we checked the first $987$ letters numerically).

\begin{remark}[Back to the channel of Lecture~11]
With $b=1$ and $a=0$, the Fibonacci word is a binary word with no two adjacent $1$'s, which is exactly
the jamming condition for discs in the narrow channel of Section~\ref{L11:sec:channel}. It is a
jammed configuration that is neither periodic nor random. Its fraction of $1$'s is
$\varphi^{-2}\approx0.382$. This lies strictly between the typical value
$\frac12-\frac{\sqrt5}{10}\approx0.276$ of Exercise~\ref{L11:ex:typical} and the maximum $\frac12$. Its
length per contact is $2\varphi^{-2}+L\varphi^{-1}$. Every finite block of it occurs in it with bounded
gaps, so it is ``as ordered as an aperiodic jammed state can be''. The count $F_{\ell+2}$ of all such
words and the lengths $\abs{\sigma^n(a)}=F_{n+2}$ obey the same Fibonacci recursion.
\end{remark}

%======================================================================
\section{Cut and project}
\label{ERG:sec:cutproject}

Here is the Fibonacci chain as a geometric object. In $\R^2$ take the line $E_\parallel$ of slope
$\varphi^{-1}$ (direction $(\varphi,1)$), the perpendicular line $E_\perp$, and the strip
$S=E_\parallel+W$, where the \emph{window} $W$ is the projection of the unit square to $E_\perp$, shifted
generically. Select the lattice points of $\Z^2$ in the strip and project them to $E_\parallel$.

\begin{figure}[ht]
\centering
\begin{tikzpicture}[scale=0.78]
  \begin{scope}
    \clip (-0.6,-2.4) rectangle (10.4,5.6);
    \fill[addcol!10] (-0.6,{(-0.6-1.3)/1.618034}) -- (10.4,{(10.4-1.3)/1.618034}) -- (10.4,{(10.4+1.318034)/1.618034}) -- (-0.6,{(-0.6+1.318034)/1.618034}) -- cycle;
    \draw[addcol!60] (-0.6,{(-0.6-1.3)/1.618034}) -- (10.4,{(10.4-1.3)/1.618034});
    \draw[addcol!60] (-0.6,{(-0.6+1.318034)/1.618034}) -- (10.4,{(10.4+1.318034)/1.618034});
  \end{scope}
  \foreach \m in {0,...,8} \foreach \k in {0,...,5} \fill[gray!60] (\m,\k) circle (0.04);
  \draw[thick,addcol] (0.667,-2.188) -- (10.18,3.691) node[right,font=\small]{$E_\parallel$};
\draw[ERGproj] (0,0)--(1.163,-1.881); \fill[addcol] (1.163,-1.881) circle (0.06);
\draw[ERGproj] (1,0)--(1.886,-1.434); \fill[addcol] (1.886,-1.434) circle (0.06);
\draw[ERGproj] (1,1)--(2.334,-1.158); \fill[addcol] (2.334,-1.158) circle (0.06);
\draw[ERGproj] (2,1)--(3.057,-0.711); \fill[addcol] (3.057,-0.711) circle (0.06);
\draw[ERGproj] (2,2)--(3.504,-0.434); \fill[addcol] (3.504,-0.434) circle (0.06);
\draw[ERGproj] (3,2)--(4.228,0.013); \fill[addcol] (4.228,0.013) circle (0.06);
\draw[ERGproj] (4,2)--(4.952,0.460); \fill[addcol] (4.952,0.460) circle (0.06);
\draw[ERGproj] (4,3)--(5.399,0.737); \fill[addcol] (5.399,0.737) circle (0.06);
\draw[ERGproj] (5,3)--(6.122,1.184); \fill[addcol] (6.122,1.184) circle (0.06);
\draw[ERGproj] (6,3)--(6.846,1.631); \fill[addcol] (6.846,1.631) circle (0.06);
\draw[ERGproj] (6,4)--(7.293,1.907); \fill[addcol] (7.293,1.907) circle (0.06);
\draw[ERGproj] (7,4)--(8.017,2.355); \fill[addcol] (8.017,2.355) circle (0.06);
\draw[ERGproj] (7,5)--(8.464,2.631); \fill[addcol] (8.464,2.631) circle (0.06);
\draw[ERGproj] (8,5)--(9.188,3.078); \fill[addcol] (9.188,3.078) circle (0.06);
\foreach \x/\y/\l in {1.672/-1.896/L, 2.257/-1.534/S, 2.843/-1.172/L, 3.428/-0.811/S, 4.013/-0.449/L,
  4.737/-0.002/L, 5.322/0.360/S, 5.908/0.722/L, 6.631/1.169/L, 7.217/1.531/S, 7.802/1.893/L,
  8.388/2.255/S, 8.973/2.617/L} \node[font=\scriptsize] at (\x,\y) {$\l$};
\draw[thick] (0,0)--(1,0)--(1,1)--(2,1)--(2,2)--(4,2)--(4,3)--(6,3)--(6,4)--(7,4)--(7,5)--(8,5);
\foreach \p in {(0,0),(1,0),(1,1),(2,1),(2,2),(3,2),(4,2),(4,3),(5,3),(6,3),(6,4),(7,4),(7,5),(8,5)} \fill \p circle (0.07);
  \draw[-{Stealth}] (9.6,0.2) -- ++(-0.4,0.647) node[left,font=\small]{$E_\perp$};
  \node[font=\small,addcol,anchor=west] at (-0.5,5.3) {strip $S=E_\parallel+W$};
\end{tikzpicture}
\caption{The Fibonacci chain by cut and project. The lattice points of $\Z^2$ in the strip (black) form a
staircase. Projected orthogonally onto $E_\parallel$ they give long ($L$) and short ($S$) intervals,
of lengths in the ratio $\varphi:1$, in the order of the Fibonacci word.}
\label{ERG:fig:cutproject}
\end{figure}

\begin{proposition}\label{ERG:prop:fibchain}
Let $w(m,k)=\varphi k-m+s$ (this is $\sqrt{1+\varphi^2}$ times the $E_\perp$-coordinate, shifted by
$s$), and select the points with $w\in(-1,\varphi]$. Each selected point $x$ is followed in the strip by
exactly one of $x+(1,0)$ (if $w(x)>0$) or $x+(0,1)$ (if $w(x)\le0$). The steps project to intervals of
lengths $\varphi/\sqrt{1+\varphi^2}$ and $1/\sqrt{1+\varphi^2}$. In the normalized coordinate
$y=(w+1)/\varphi^2\in(0,1]$, every step acts as $y\mapsto y-\varphi^{-2}\pmod1$, and the step is short
exactly when $y\le\varphi^{-2}$. So the sequence of tiles is a mechanical word of slope $\varphi^{-2}$:
a Fibonacci word (for $s$ chosen suitably, exactly $f$).
\end{proposition}

\begin{proof}
If $w(x)>0$ then $w(x+(1,0))=w(x)-1\in(-1,\varphi-1]$ is in the window, while $w(x+(0,1))=w(x)+\varphi>\varphi$
is not. If $w(x)\le0$ the roles are reversed, since $w(x)-1\le-1$ and $w(x)+\varphi\in(\varphi-1,\varphi]$.
A step $(1,0)$ changes $y$ by $-1/\varphi^2$, and a step $(0,1)$ by $\varphi/\varphi^2=1-\varphi^{-2}$,
and $w>0$ means $y>\varphi^{-2}$. It remains to see that the staircase through one selected point
contains all of them. Every selected point lies on such a staircase (run the rule backwards as well), and
distinct staircases are disjoint. By Theorem~\ref{ERG:thm:weyl} a staircase has a fraction
$\varphi^{-2}$ of short steps, so it has $1/(\varphi^{-1}\cdot\frac{\varphi}{\sqrt{1+\varphi^2}}+
\varphi^{-2}\frac{1}{\sqrt{1+\varphi^2}})=\varphi^2/\sqrt{1+\varphi^2}$ points per unit length of
$E_\parallel$. This is already the density of all lattice points in the strip, which is (strip width)/(covolume) $=\varphi^2/\sqrt{1+\varphi^2}$. So there is only one staircase.
\end{proof}

This is the prototype of a \emph{model set}.

\begin{definition}
A \emph{cut and project scheme} is a lattice $\mathcal L\subset E_\parallel\times E_\perp$ such that the
projection $\pi_\parallel$ is injective on $\mathcal L$ and $\pi_\perp(\mathcal L)$ is dense in
$E_\perp$. For a \emph{window} $W\subset E_\perp$, relatively compact with nonempty interior, the
\emph{model set} is $\Lambda(W)=\{\pi_\parallel(x):x\in\mathcal L,\ \pi_\perp(x)\in W\}$. It is
\emph{regular} if $\partial W$ has measure zero.
\end{definition}

A model set is a Delone set (uniformly discrete and relatively dense), so it is a packing in the sense of
Lecture~1 and has a density. It is in general not periodic. It is a \emph{Meyer set}:
$\Lambda-\Lambda\subseteq\Lambda+F$ for a finite set $F$ (Meyer 1972). Its density is
$\vol(W)/\covol(\mathcal L)$, by Weyl equidistribution of $\pi_\perp(\mathcal L)$ in $W$. Most importantly
for physics, \emph{regular model sets have pure point diffraction}: the Fourier transform of the
autocorrelation of the Dirac comb $\sum_{x\in\Lambda}\delta_x$ is a sum of point masses, supported on the
projection of the dual lattice $\mathcal L^*$ (Hof 1995 for Euclidean internal spaces; Schlottmann 2000
in general). The proof is a Poisson summation formula (Interlude~B) for $\mathcal L$, cut down by the
window.

\paragraph{Penrose tilings and Shechtman's crystals.} De Bruijn (1981) showed that the rhombic Penrose
tilings (Penrose 1974) are dual to \emph{pentagrids}: superpositions of five families of parallel lines
at angles $2\pi j/5$. Equivalently, their vertex sets are obtained by cut and project from $\Z^5$, with
$E_\parallel$ the plane on which the cyclic permutation of coordinates acts as rotation by $2\pi/5$. Since
the vertices lie in finitely many layers along $(1,1,1,1,1)$, the minimal setting is the
four-dimensional root lattice $A_4$, with a window made of four pentagons in a two-dimensional internal
space (Baake--Kramer--Schlottmann--Zeidler 1990). Figure~\ref{ERG:fig:penrose} was drawn from a
pentagrid. In 1984 Shechtman, Blech, Gratias and Cahn reported a rapidly
solidified Al--Mn alloy whose electron diffraction pattern had sharp spots and icosahedral symmetry. Such
symmetry is impossible for a crystal: a rotation of order $5$ preserving a lattice in $\R^2$ or $\R^3$
would have integer trace, but its trace is $2\cos\frac{2\pi}5=\varphi^{-1}$, or $1+\varphi^{-1}$. The
discovery earned Shechtman the 2011 Nobel Prize in Chemistry. Model sets with icosahedral symmetry are
the standard mathematical models of such quasicrystals.

\begin{figure}[ht]
\centering
\begin{tikzpicture}[scale=0.52]
\begin{scope}
\clip (0,0) circle (4.9);
\filldraw[ERGthick] (-5.545,0.951)--(-4.545,0.951)--(-4.236,1.902)--(-5.236,1.902)--cycle;
\filldraw[ERGthick] (-3.427,-2.490)--(-2.427,-2.490)--(-2.118,-1.539)--(-3.118,-1.539)--cycle;
\filldraw[ERGthick] (-3.427,-0.588)--(-2.427,-0.588)--(-2.118,0.363)--(-3.118,0.363)--cycle;
\filldraw[ERGthick] (-3.427,2.490)--(-2.427,2.490)--(-2.118,3.441)--(-3.118,3.441)--cycle;
\filldraw[ERGthick] (-0.809,-5.568)--(0.191,-5.568)--(0.500,-4.617)--(-0.500,-4.617)--cycle;
\filldraw[ERGthick] (-1.618,-3.078)--(-0.618,-3.078)--(-0.309,-2.127)--(-1.309,-2.127)--cycle;
\filldraw[ERGthick] (-0.809,-0.588)--(0.191,-0.588)--(0.500,0.363)--(-0.500,0.363)--cycle;
\filldraw[ERGthick] (-1.618,1.902)--(-0.618,1.902)--(-0.309,2.853)--(-1.309,2.853)--cycle;
\filldraw[ERGthick] (-0.809,4.392)--(0.191,4.392)--(0.500,5.343)--(-0.500,5.343)--cycle;
\filldraw[ERGthick] (1.309,-4.029)--(2.309,-4.029)--(2.618,-3.078)--(1.618,-3.078)--cycle;
\filldraw[ERGthick] (1.309,-2.127)--(2.309,-2.127)--(2.618,-1.176)--(1.618,-1.176)--cycle;
\filldraw[ERGthick] (1.309,0.951)--(2.309,0.951)--(2.618,1.902)--(1.618,1.902)--cycle;
\filldraw[ERGthick] (2.118,3.441)--(3.118,3.441)--(3.427,4.392)--(2.427,4.392)--cycle;
\filldraw[ERGthick] (3.927,-2.127)--(4.927,-2.127)--(5.236,-1.176)--(4.236,-1.176)--cycle;
\filldraw[ERGthick] (3.927,0.951)--(4.927,0.951)--(5.236,1.902)--(4.236,1.902)--cycle;
\filldraw[ERGthin] (-2.618,-3.078)--(-1.618,-3.078)--(-2.427,-2.490)--(-3.427,-2.490)--cycle;
\filldraw[ERGthin] (-2.618,1.902)--(-1.618,1.902)--(-2.427,2.490)--(-3.427,2.490)--cycle;
\filldraw[ERGthin] (-0.809,-3.665)--(0.191,-3.665)--(-0.618,-3.078)--(-1.618,-3.078)--cycle;
\filldraw[ERGthin] (-0.500,0.363)--(0.500,0.363)--(-0.309,0.951)--(-1.309,0.951)--cycle;
\filldraw[ERGthin] (2.118,-4.617)--(3.118,-4.617)--(2.309,-4.029)--(1.309,-4.029)--cycle;
\filldraw[ERGthin] (2.118,0.363)--(3.118,0.363)--(2.309,0.951)--(1.309,0.951)--cycle;
\filldraw[ERGthin] (4.236,-1.176)--(5.236,-1.176)--(4.427,-0.588)--(3.427,-0.588)--cycle;
\filldraw[ERGthin] (-2.618,1.902)--(-1.618,1.902)--(-2.427,1.314)--(-3.427,1.314)--cycle;
\filldraw[ERGthin] (-2.618,-3.078)--(-1.618,-3.078)--(-2.427,-3.665)--(-3.427,-3.665)--cycle;
\filldraw[ERGthin] (-0.500,3.441)--(0.500,3.441)--(-0.309,2.853)--(-1.309,2.853)--cycle;
\filldraw[ERGthin] (-0.500,-1.539)--(0.500,-1.539)--(-0.309,-2.127)--(-1.309,-2.127)--cycle;
\filldraw[ERGthin] (2.118,3.441)--(3.118,3.441)--(2.309,2.853)--(1.309,2.853)--cycle;
\filldraw[ERGthin] (2.427,-0.588)--(3.427,-0.588)--(2.618,-1.176)--(1.618,-1.176)--cycle;
\filldraw[ERGthin] (4.236,-0.000)--(5.236,-0.000)--(4.427,-0.588)--(3.427,-0.588)--cycle;
\filldraw[ERGthick] (-3.427,4.392)--(-2.427,4.392)--(-2.118,3.441)--(-3.118,3.441)--cycle;
\filldraw[ERGthick] (-3.427,1.314)--(-2.427,1.314)--(-2.118,0.363)--(-3.118,0.363)--cycle;
\filldraw[ERGthick] (-3.427,-0.588)--(-2.427,-0.588)--(-2.118,-1.539)--(-3.118,-1.539)--cycle;
\filldraw[ERGthick] (-3.427,-3.665)--(-2.427,-3.665)--(-2.118,-4.617)--(-3.118,-4.617)--cycle;
\filldraw[ERGthick] (-0.809,4.392)--(0.191,4.392)--(0.500,3.441)--(-0.500,3.441)--cycle;
\filldraw[ERGthick] (-1.618,1.902)--(-0.618,1.902)--(-0.309,0.951)--(-1.309,0.951)--cycle;
\filldraw[ERGthick] (-0.809,-0.588)--(0.191,-0.588)--(0.500,-1.539)--(-0.500,-1.539)--cycle;
\filldraw[ERGthick] (-0.809,-3.665)--(0.191,-3.665)--(0.500,-4.617)--(-0.500,-4.617)--cycle;
\filldraw[ERGthick] (1.309,2.853)--(2.309,2.853)--(2.618,1.902)--(1.618,1.902)--cycle;
\filldraw[ERGthick] (2.118,0.363)--(3.118,0.363)--(3.427,-0.588)--(2.427,-0.588)--cycle;
\filldraw[ERGthick] (1.309,-2.127)--(2.309,-2.127)--(2.618,-3.078)--(1.618,-3.078)--cycle;
\filldraw[ERGthick] (3.927,2.853)--(4.927,2.853)--(5.236,1.902)--(4.236,1.902)--cycle;
\filldraw[ERGthick] (3.927,0.951)--(4.927,0.951)--(5.236,-0.000)--(4.236,-0.000)--cycle;
\filldraw[ERGthick] (3.118,-4.617)--(3.427,-3.665)--(2.618,-3.078)--(2.309,-4.029)--cycle;
\filldraw[ERGthick] (0.191,-3.665)--(0.500,-2.714)--(-0.309,-2.127)--(-0.618,-3.078)--cycle;
\filldraw[ERGthick] (-1.618,-3.078)--(-1.309,-2.127)--(-2.118,-1.539)--(-2.427,-2.490)--cycle;
\filldraw[ERGthick] (-4.545,-2.127)--(-4.236,-1.176)--(-5.045,-0.588)--(-5.354,-1.539)--cycle;
\filldraw[ERGthick] (3.927,-2.127)--(4.236,-1.176)--(3.427,-0.588)--(3.118,-1.539)--cycle;
\filldraw[ERGthick] (1.309,-2.127)--(1.618,-1.176)--(0.809,-0.588)--(0.500,-1.539)--cycle;
\filldraw[ERGthick] (-0.809,-0.588)--(-0.500,0.363)--(-1.309,0.951)--(-1.618,0.000)--cycle;
\filldraw[ERGthick] (-3.736,0.363)--(-3.427,1.314)--(-4.236,1.902)--(-4.545,0.951)--cycle;
\filldraw[ERGthick] (3.118,0.363)--(3.427,1.314)--(2.618,1.902)--(2.309,0.951)--cycle;
\filldraw[ERGthick] (1.000,1.902)--(1.309,2.853)--(0.500,3.441)--(0.191,2.490)--cycle;
\filldraw[ERGthick] (-1.618,1.902)--(-1.309,2.853)--(-2.118,3.441)--(-2.427,2.490)--cycle;
\filldraw[ERGthick] (3.927,2.853)--(4.236,3.804)--(3.427,4.392)--(3.118,3.441)--cycle;
\filldraw[ERGthick] (2.118,3.441)--(2.427,4.392)--(1.618,4.980)--(1.309,4.029)--cycle;
\filldraw[ERGthin] (1.000,-3.078)--(1.309,-2.127)--(0.500,-2.714)--(0.191,-3.665)--cycle;
\filldraw[ERGthin] (-3.736,-1.539)--(-3.427,-0.588)--(-4.236,-1.176)--(-4.545,-2.127)--cycle;
\filldraw[ERGthin] (3.118,-1.539)--(3.427,-0.588)--(2.618,-1.176)--(2.309,-2.127)--cycle;
\filldraw[ERGthin] (-1.618,0.000)--(-1.309,0.951)--(-2.118,0.363)--(-2.427,-0.588)--cycle;
\filldraw[ERGthin] (3.927,0.951)--(4.236,1.902)--(3.427,1.314)--(3.118,0.363)--cycle;
\filldraw[ERGthin] (0.191,2.490)--(0.500,3.441)--(-0.309,2.853)--(-0.618,1.902)--cycle;
\filldraw[ERGthin] (1.309,4.029)--(1.618,4.980)--(0.809,4.392)--(0.500,3.441)--cycle;
\filldraw[ERGthin] (-3.736,-1.539)--(-3.427,-0.588)--(-3.118,-1.539)--(-3.427,-2.490)--cycle;
\filldraw[ERGthin] (1.000,-3.078)--(1.309,-2.127)--(1.618,-3.078)--(1.309,-4.029)--cycle;
\filldraw[ERGthin] (-3.736,0.363)--(-3.427,1.314)--(-3.118,0.363)--(-3.427,-0.588)--cycle;
\filldraw[ERGthin] (0.191,-0.588)--(0.500,0.363)--(0.809,-0.588)--(0.500,-1.539)--cycle;
\filldraw[ERGthin] (-3.736,3.441)--(-3.427,4.392)--(-3.118,3.441)--(-3.427,2.490)--cycle;
\filldraw[ERGthin] (1.000,1.902)--(1.309,2.853)--(1.618,1.902)--(1.309,0.951)--cycle;
\filldraw[ERGthin] (0.191,4.392)--(0.500,5.343)--(0.809,4.392)--(0.500,3.441)--cycle;
\filldraw[ERGthick] (3.927,-2.127)--(3.118,-1.539)--(2.309,-2.127)--(3.118,-2.714)--cycle;
\filldraw[ERGthick] (2.118,-4.617)--(1.309,-4.029)--(0.500,-4.617)--(1.309,-5.204)--cycle;
\filldraw[ERGthick] (4.236,1.902)--(3.427,2.490)--(2.618,1.902)--(3.427,1.314)--cycle;
\filldraw[ERGthick] (2.118,0.363)--(1.309,0.951)--(0.500,0.363)--(1.309,-0.225)--cycle;
\filldraw[ERGthick] (1.309,-2.127)--(0.500,-1.539)--(-0.309,-2.127)--(0.500,-2.714)--cycle;
\filldraw[ERGthick] (-0.500,-4.617)--(-1.309,-4.029)--(-2.118,-4.617)--(-1.309,-5.204)--cycle;
\filldraw[ERGthick] (2.118,3.441)--(1.309,4.029)--(0.500,3.441)--(1.309,2.853)--cycle;
\filldraw[ERGthick] (1.000,1.902)--(0.191,2.490)--(-0.618,1.902)--(0.191,1.314)--cycle;
\filldraw[ERGthick] (-0.809,-0.588)--(-1.618,0.000)--(-2.427,-0.588)--(-1.618,-1.176)--cycle;
\filldraw[ERGthick] (-2.618,-3.078)--(-3.427,-2.490)--(-4.236,-3.078)--(-3.427,-3.665)--cycle;
\filldraw[ERGthick] (-0.809,4.392)--(-1.618,4.980)--(-2.427,4.392)--(-1.618,3.804)--cycle;
\filldraw[ERGthick] (-2.618,1.902)--(-3.427,2.490)--(-4.236,1.902)--(-3.427,1.314)--cycle;
\filldraw[ERGthick] (-3.427,-0.588)--(-4.236,0.000)--(-5.045,-0.588)--(-4.236,-1.176)--cycle;
\filldraw[ERGthin] (3.118,-2.714)--(2.309,-2.127)--(2.618,-3.078)--(3.427,-3.665)--cycle;
\filldraw[ERGthin] (3.927,2.853)--(3.118,3.441)--(3.427,2.490)--(4.236,1.902)--cycle;
\filldraw[ERGthin] (1.309,-0.225)--(0.500,0.363)--(0.809,-0.588)--(1.618,-1.176)--cycle;
\filldraw[ERGthin] (-0.809,-3.665)--(-1.618,-3.078)--(-1.309,-4.029)--(-0.500,-4.617)--cycle;
\filldraw[ERGthin] (0.191,1.314)--(-0.618,1.902)--(-0.309,0.951)--(0.500,0.363)--cycle;
\filldraw[ERGthin] (-1.618,-1.176)--(-2.427,-0.588)--(-2.118,-1.539)--(-1.309,-2.127)--cycle;
\filldraw[ERGthin] (-1.618,3.804)--(-2.427,4.392)--(-2.118,3.441)--(-1.309,2.853)--cycle;
\filldraw[ERGthin] (-3.736,0.363)--(-4.545,0.951)--(-4.236,0.000)--(-3.427,-0.588)--cycle;
\filldraw[ERGthick] (3.118,3.441)--(2.309,2.853)--(2.618,1.902)--(3.427,2.490)--cycle;
\filldraw[ERGthick] (3.927,0.951)--(3.118,0.363)--(3.427,-0.588)--(4.236,-0.000)--cycle;
\filldraw[ERGthick] (-0.809,4.392)--(-1.618,3.804)--(-1.309,2.853)--(-0.500,3.441)--cycle;
\filldraw[ERGthick] (1.000,1.902)--(0.191,1.314)--(0.500,0.363)--(1.309,0.951)--cycle;
\filldraw[ERGthick] (2.118,0.363)--(1.309,-0.225)--(1.618,-1.176)--(2.427,-0.588)--cycle;
\filldraw[ERGthick] (3.927,-2.127)--(3.118,-2.714)--(3.427,-3.665)--(4.236,-3.078)--cycle;
\filldraw[ERGthick] (-3.736,3.441)--(-4.545,2.853)--(-4.236,1.902)--(-3.427,2.490)--cycle;
\filldraw[ERGthick] (-1.618,1.902)--(-2.427,1.314)--(-2.118,0.363)--(-1.309,0.951)--cycle;
\filldraw[ERGthick] (-0.809,-0.588)--(-1.618,-1.176)--(-1.309,-2.127)--(-0.500,-1.539)--cycle;
\filldraw[ERGthick] (1.000,-3.078)--(0.191,-3.665)--(0.500,-4.617)--(1.309,-4.029)--cycle;
\filldraw[ERGthick] (-4.545,0.951)--(-5.354,0.363)--(-5.045,-0.588)--(-4.236,0.000)--cycle;
\filldraw[ERGthick] (-3.736,-1.539)--(-4.545,-2.127)--(-4.236,-3.078)--(-3.427,-2.490)--cycle;
\filldraw[ERGthick] (-1.618,-3.078)--(-2.427,-3.665)--(-2.118,-4.617)--(-1.309,-4.029)--cycle;
\end{scope}
\end{tikzpicture}
\caption{A patch of a rhombic Penrose tiling, computed from a generic pentagrid (de Bruijn's method).
Thick rhombi ($72^\circ$) blue, thin rhombi ($36^\circ$) orange.}
\label{ERG:fig:penrose}
\end{figure}

%======================================================================
\section{Folding Dynkin diagrams: quasicrystals from \texorpdfstring{$A_4$, $D_6$, $E_8$}{A4, D6, E8}}
\label{ERG:sec:fold}

Where do the icosahedral windows and lattices come from? Return to Lecture~4. Dropping the
crystallographic axiom~4) from Definition~\ref{L04a:def:rootsystem} allows more root systems. The
exercise after that definition asked for the decagon and the icosidodecahedron. Coxeter's classification
of finite reflection groups (1934) adds exactly three types to the list of Lecture~5: the dihedral groups
$I_2(m)$ ($m\neq2,3,4,6$), and two exceptional groups $H_3$ and $H_4$, with Coxeter graphs
\[
  H_2=I_2(5):\ \circ\overset{5}{\text{---}}\circ,\qquad
  H_3:\ \circ\text{---}\circ\overset{5}{\text{---}}\circ,\qquad
  H_4:\ \circ\text{---}\circ\text{---}\circ\overset{5}{\text{---}}\circ .
\]
These are the symmetry groups of the regular pentagon (order $10$), the icosahedron (order $120$) and the
$600$-cell (order $14400$). An edge labelled $5$ means an angle of $\pi-\pi/5$ between simple roots, so
$\cos\frac\pi5=\frac\tau2$ enters, and the Cartan numbers are irrational. These groups preserve no lattice
in their own dimension. But they do preserve a lattice of \emph{twice} the dimension, and that is where
the quasicrystals come from.

The construction is a \emph{folding} of a simply laced Dynkin diagram (Shcherbak 1988; Moody--Patera
1993). Pair up the $2r$ nodes of $\Gamma=A_4,D_6,E_8$ as $(\alpha_i,\beta_i)$, $1\le i\le r$ (so $r=2,3,4$), as in
Figure~\ref{ERG:fig:fold}. The rule is simple. Each plain edge $i$---$j$ of $H_r$ gives the two edges
$\alpha_i$---$\alpha_j$ and $\beta_i$---$\beta_j$. Each $5$-edge $i$---$j$ gives the three edges
$\alpha_i$---$\beta_j$, $\beta_i$---$\alpha_j$ and $\beta_i$---$\beta_j$. For $H_4$ this produces the
tree with arms of lengths $4,2,1$ from the node $\beta_3$, which is $E_8$. For $H_3$ it produces $D_6$, and
for $H_2$ the chain $\alpha_1$---$\beta_2$---$\beta_1$---$\alpha_2$, which is $A_4$.

\begin{figure}[ht]
\centering
\begin{tikzpicture}[scale=0.86]
  % E8 -> H4
  \begin{scope}
    \foreach \i in {1,...,4} {\node[ERGv] (a\i) at (\i-1,1) {}; \node[font=\scriptsize,above=2pt] at (\i-1,1) {$\alpha_{\i}$};
                              \node[ERGw] (b\i) at (\i-1,0) {}; \node[font=\scriptsize,below=2pt] at (\i-1,0) {$\beta_{\i}$};
                              \draw[gray,dashed] (a\i)--(b\i);}
    \draw (a1)--(a2)--(a3); \draw (b1)--(b2)--(b3)--(b4); \draw (a3)--(b4); \draw (a4)--(b3);
    \node[font=\small] at (1.5,-0.85) {$E_8$};
    \draw[-{Stealth},thick] (3.5,0.5)--(4.3,0.5);
    \foreach \i in {1,...,4} {\node[ERGv] (h\i) at (4.6+0.7*\i-0.7,0.5) {}; \node[font=\scriptsize,above=2pt] at (4.6+0.7*\i-0.7,0.5) {$a_{\i}$};}
    \draw (h1)--(h2)--(h3)--node[below,font=\scriptsize]{$5$}(h4);
    \node[font=\small] at (5.65,-0.85) {$H_4$};
  \end{scope}
  % D6 -> H3
  \begin{scope}[xshift=7.6cm]
    \foreach \i in {1,...,3} {\node[ERGv] (a\i) at (\i-1,1) {}; \node[font=\scriptsize,above=2pt] at (\i-1,1) {$\alpha_{\i}$};
                              \node[ERGw] (b\i) at (\i-1,0) {}; \node[font=\scriptsize,below=2pt] at (\i-1,0) {$\beta_{\i}$};
                              \draw[gray,dashed] (a\i)--(b\i);}
    \draw (a1)--(a2); \draw (b1)--(b2)--(b3); \draw (a2)--(b3); \draw (a3)--(b2);
    \node[font=\small] at (1,-0.85) {$D_6$};
    \draw[-{Stealth},thick] (2.5,0.5)--(3.2,0.5);
    \foreach \i in {1,...,3} {\node[ERGv] (h\i) at (3.5+0.7*\i-0.7,0.5) {}; \node[font=\scriptsize,above=2pt] at (3.5+0.7*\i-0.7,0.5) {$a_{\i}$};}
    \draw (h1)--(h2)--node[below,font=\scriptsize]{$5$}(h3);
    \node[font=\small] at (4.2,-0.85) {$H_3$};
  \end{scope}
  % A4 -> H2
  \begin{scope}[xshift=13.3cm]
    \foreach \i in {1,2} {\node[ERGv] (a\i) at (\i-1,1) {}; \node[font=\scriptsize,above=2pt] at (\i-1,1) {$\alpha_{\i}$};
                          \node[ERGw] (b\i) at (\i-1,0) {}; \node[font=\scriptsize,below=2pt] at (\i-1,0) {$\beta_{\i}$};
                          \draw[gray,dashed] (a\i)--(b\i);}
    \draw (b1)--(b2); \draw (a1)--(b2); \draw (a2)--(b1);
    \node[font=\small] at (0.5,-0.85) {$A_4$};
    \draw[-{Stealth},thick] (1.4,0.5)--(2.0,0.5);
    \node[ERGv] (h1) at (2.3,0.5) {}; \node[ERGv] (h2) at (3.0,0.5) {};
    \node[font=\scriptsize,above=2pt] at (2.3,0.5) {$a_1$}; \node[font=\scriptsize,above=2pt] at (3.0,0.5) {$a_2$};
    \draw (h1)--node[below,font=\scriptsize]{$5$}(h2);
    \node[font=\small] at (2.65,-0.85) {$H_2$};
  \end{scope}
\end{tikzpicture}
\caption{Folding $E_8\to H_4$, $D_6\to H_3$, $A_4\to H_2$. The dashed lines pair each node $\alpha_i$
with its partner $\beta_i$ (shaded). These are orthogonal, and under the projection
$\beta_i\mapsto\tau a_i$, $\alpha_i\mapsto a_i$. A plain edge of $H_r$ comes from two parallel edges, and a
$5$-edge from a crossing pair plus the bottom edge. The $E_8$ diagram here is the chain
$\alpha_1\alpha_2\alpha_3\beta_4\beta_3\beta_2\beta_1$ with $\alpha_4$ attached to $\beta_3$.}
\label{ERG:fig:fold}
\end{figure}

\begin{theorem}[Folding]\label{ERG:thm:fold}
Let $\Gamma\in\{A_4,D_6,E_8\}$ with simple roots labelled as in Figure~\ref{ERG:fig:fold}, $L$ its root
lattice in $V=\R^{2r}$, and define $M\colon L\to L$ on the simple roots by
\[
  M\alpha_i=\beta_i,\qquad M\beta_i=\alpha_i+\beta_i .
\]
\begin{enumerate}[nosep]
\item $M$ is self-adjoint and $M^2=M+1$. Hence $V=E_\tau\oplus E_{\tau'}$ is an orthogonal sum of the
  $\tau$- and $\tau'$-eigenspaces of $M$, each of dimension $r$, and $L$ is a module over
  $\Z[\tau]$ (with $\tau$ acting as $M$).
\item Let $P$ be the orthogonal projection onto $E_\tau$ and $a_i=P\alpha_i$. Then $P\beta_i=\tau a_i$, and
  $a_1,\dots,a_r$ are the simple roots of $H_r$, all of squared length $1-\frac1{\sqrt5}$.
\item The products $w_i=s_{\alpha_i}s_{\beta_i}$ preserve $E_\tau$ and act on it as the reflections
  $s_{a_i}$. They generate a subgroup of the Weyl group $W(\Gamma)$ isomorphic to $W(H_r)$.
\item $P$ is injective on $L$, and $P(L)=\Z[\tau]a_1+\dots+\Z[\tau]a_r$ is dense in $E_\tau$.
\end{enumerate}
\end{theorem}

\begin{proof}
(1) $M$ is self-adjoint if and only if $\ip{\beta_i}{\alpha_j}=\ip{\alpha_i}{\beta_j}$ and
$\ip{\beta_i}{\beta_j}=\ip{\alpha_i}{\alpha_j}+\ip{\alpha_i}{\beta_j}$ for all $i,j$. Both identities can
be read off the edge rule. Each side of the first is $-1$ exactly for a $5$-edge. For the second, take
$i=j$ ($2=2+0$), a plain edge ($-1=-1+0$), a $5$-edge ($-1=0-1$), and no edge ($0=0$). Clearly
$M^2\alpha_i=M\beta_i=\alpha_i+\beta_i=(M+1)\alpha_i$, and similarly for $\beta_i$. A self-adjoint
operator with $M^2=M+1$ has eigenvalues among $\tau,\tau'$ and orthogonal eigenspaces. Since $M$ has an
integer matrix, the two eigenvalues, being Galois conjugate, have equal multiplicity.

(2) $P=(M-\tau')/\sqrt5$ commutes with $M$, so $P\beta_i=PM\alpha_i=MP\alpha_i=\tau a_i$. Then
$\ip{a_i}{a_j}=\ip{P\alpha_i}{\alpha_j}=\big(\ip{\beta_i}{\alpha_j}-\tau'\ip{\alpha_i}{\alpha_j}\big)/\sqrt5$.
This gives $-2\tau'/\sqrt5=1-\frac1{\sqrt5}=:x$ for $i=j$, $\tau'/\sqrt5=-\frac x2$ for a plain edge, so
the angle is $2\pi/3$, and $-1/\sqrt5=-x\frac\tau2$ for a $5$-edge, so the angle is $\pi-\pi/5$. Non-adjacent
pairs give $0$.

(3) The roots $\alpha_i,\beta_i$ are orthogonal, so $w_i=-1$ on $U_i=\operatorname{span}(\alpha_i,\beta_i)$
and $w_i=1$ on $U_i^\perp$. The plane $U_i$ is $M$-invariant, hence so is $U_i^\perp$, so $w_i$ commutes
with $M$ and preserves $E_\tau$. On $E_\tau$ its $(-1)$-eigenspace is $U_i\cap E_\tau=\R a_i$. So
$w_i|_{E_\tau}=s_{a_i}$, and these generate $W(H_r)$. If some $w\in\langle w_i\rangle$ is trivial on
$E_\tau$, then $w-1$ is a rational matrix vanishing on $E_\tau$, hence also on its Galois conjugate
$E_{\tau'}$, so $w=1$.

(4) A nonzero $v\in L$ with $Pv=0$ would satisfy $Mv=\tau'v$, impossible for a rational vector $v$ and an
integer matrix $M$. The image is the $\Z[\tau]$-span of the $a_i$, which is dense because $\Z[\tau]$ is
dense in $\R$.
\end{proof}

A direct computation (we did it in python3) shows that $P$ maps the roots of $\Gamma$ onto
$\Delta(H_r)\cup\tau\Delta(H_r)$:
\[
  A_4:\ 20=10+10,\qquad D_6:\ 60=30+30,\qquad E_8:\ 240=120+120 .
\]
Here $\Delta(H_r)$ is the root system of $H_r$: the vertices of a decagon, of an icosidodecahedron, and
of a $600$-cell. The $120$ unit quaternions forming the $600$-cell are the \emph{icosians}; their
$\Z[\tau]$-span, the icosian ring, is a $\Z[\tau]$-order in the quaternions, and Theorem~\ref{ERG:thm:fold}(1) is a shadow of the fact that $E_8$ is the
icosian ring, with a suitably modified norm, viewed as a lattice over $\Z$ (Conway--Sloane, Ch.~8, \S2; Moody--Patera 1993). The planar
case is drawn in Figure~\ref{ERG:fig:a4}.

\begin{figure}[ht]
\centering
\begin{tikzpicture}[scale=1.25]
  \foreach \k in {0,...,9} {\draw[gray!40] (0,0)--({36*\k}:1.618);}
  \draw[addcol] (0:1) \foreach \k in {1,...,10} {-- ({36*\k}:1)};
  \draw[flagcol] (0:1.618) \foreach \k in {1,...,10} {-- ({36*\k}:1.618)};
  \foreach \k in {0,...,9} {\fill[addcol] ({36*\k}:1) circle (0.045); \fill[flagcol] ({36*\k}:1.618) circle (0.045);}
  \draw[-{Stealth},thick] (0,0)--(0:1) node[below right=-1pt,font=\scriptsize]{$a_1$};
  \draw[-{Stealth},thick] (0,0)--(144:1) node[below left=-1pt,font=\scriptsize]{$a_2$};
  \node[font=\scriptsize,flagcol] at (0:2.05) {$\tau a_1$};
  \node[font=\small,align=left,anchor=west] at (2.4,0.3) {$P(\alpha)$ for the $20$ roots of $A_4$:};
  \node[font=\small,align=left,anchor=west] at (2.4,-0.15) {\textcolor{addcol}{inner decagon} $=\Delta(H_2)$,};
  \node[font=\small,align=left,anchor=west] at (2.4,-0.6) {\textcolor{flagcol}{outer decagon} $=\tau\,\Delta(H_2)$.};
\end{tikzpicture}
\caption{The $20$ roots of $A_4$ projected to the plane $E_\tau$ (the Coxeter plane of $A_4$). The two
concentric decagons have radii in the ratio $\tau$.}
\label{ERG:fig:a4}
\end{figure}

\paragraph{Quasicrystals.} By Theorem~\ref{ERG:thm:fold}(4), $(L,E_\tau,E_{\tau'})$ is a cut and project
scheme whose symmetry group $W(H_r)$ acts on both factors, on $E_\tau$ as $H_r$ and on $E_{\tau'}$ as its
Galois twin. Choosing an $H_r$-invariant window in $E_{\tau'}$ gives model sets with decagonal ($A_4$),
icosahedral ($D_6$; face-centred icosahedral phases such as Al--Cu--Fe are modelled on it, while
primitive ones such as Shechtman's Al--Mn use $\Z^6$) or $H_4$ symmetry ($E_8$, the
four-dimensional quasicrystal of Elser--Sloane 1987). In this way the lattices of Lecture~6 \emph{are}
the hidden periodic structures behind the quasicrystals.

\paragraph{Affine extensions and viruses.} For a crystallographic diagram, adding the extra node of the
extended (affine) Dynkin diagram adds a translation, and the group becomes an infinite discrete affine
reflection group, as in Interlude~A's alcoves. For $H_r$, adding a translation produces a group whose
orbits are dense, because $\Z[\tau]$ is. Patera and Twarock (2002) constructed an affine extension of each of
$H_2,H_3,H_4$ (unique under their axioms), and used finitely many steps of it to generate finite
quasicrystal fragments.
Dechant, Boehm and Twarock (2013) showed that some of these extensions are exactly what one gets by
folding the \emph{affine} diagrams $\tilde A_4,\tilde D_6,\tilde E_8$ and projecting. Twarock's
motivation was biology. Many viral capsids are icosahedral. Caspar--Klug theory explains many of them,
but not, for instance, the papovaviruses, whose $72$ capsomers are all pentamers. Twarock (2004)
explained their architecture by tilings built from the icosahedral group. Keef and Twarock (2009) used
affine extensions of the icosahedral group to describe how material at different radii of a virus
(capsid, genome) is organised by one common icosahedral blueprint.

%======================================================================
\section{The space of lattices as a dynamical system}
\label{ERG:sec:lattices}

Section~\ref{ERG:sec:dioph} used the space $X_n=\SL_n(\R)/\SL_n(\Z)$ of unimodular lattices from
Lecture~4. It is not compact (Mahler's criterion characterizes its bounded sets), but it has finite
volume. Siegel's mean value theorem (Interlude~B) says that, for the normalized invariant probability
measure $\mu_n$ on $X_n$ and $f$ integrable on $\R^n$,
\[
  \int_{X_n}\sum_{v\in\Lat\setminus0}f(v)\,d\mu_n(\Lat)=\int_{\R^n}f .
\]
Take $f=\mathbf 1_B$ for a ball $B$ of radius $r$ and volume $V<2$. The average number of nonzero
lattice vectors in $B$ is $V<2$. Since nonzero lattice vectors come in pairs $\pm v$, \emph{some} lattice
has none. Its balls of radius $r/2$ form a packing of density $V2^{-n}$. This gives the Minkowski--Hlawka bound
$\dens\ge2^{1-n}$ for lattice packings, and counting only primitive vectors improves it to
$\zeta(n)2^{1-n}$. It is a non-constructive existence statement, like the saturation bound of
Lecture~1, but for lattices.

$\SL_n(\R)$ acts on $X_n$, and the dynamics of its subgroups encode arithmetic. The diagonal flow
$g_t$ of Proposition~\ref{ERG:prop:dani} is the geodesic flow on the modular surface. \emph{Unipotent}
subgroups behave much more rigidly.

\begin{theorem}[Margulis 1987, the Oppenheim conjecture]\label{ERG:thm:margulis}
Let $Q$ be a real nondegenerate indefinite quadratic form in $n\ge3$ variables that is not a multiple of a
form with rational coefficients. Then for every $\varepsilon>0$ there is $x\in\Z^n\setminus0$ with
$\abs{Q(x)}<\varepsilon$. (Dani and Margulis, \emph{Invent.\ Math.}~98 (1989), strengthened this to the
density of $Q(x)$, $x\in\Z^n$ primitive, in $\R$.)
\end{theorem}

Margulis's proof is a statement about orbits. It suffices to treat $n=3$. Write $Q=Q_0\circ g$ with
$Q_0$ a fixed form of signature $(2,1)$, and let $H=\SO(Q_0)^\circ\cong\SO(2,1)^\circ$, which is generated by
unipotent elements. Raghunathan observed that it suffices to show: \emph{every bounded $H$-orbit in $X_3$ is
compact}. Margulis proved this. A compact orbit $Hg\Z^3$ occurs only when $Q$ is a multiple of a rational
form. Otherwise the orbit is unbounded, so by Mahler's criterion some $hg\Z^3$ contains a very short vector
$hgx$, and then $\abs{Q(x)}=\abs{Q_0(hgx)}$ is small. Theorem~\ref{ERG:thm:margulis} fails for $n=2$, and the
reason is Section~\ref{ERG:sec:dioph}. For $Q(x,y)=(x-\sqrt2y)(x-\sqrt3y)$ one has
$\inf_{\Z^2\setminus0}\abs{Q}>0$ (Exercise~\ref{ERG:ex:oppenheim2}), because $\sqrt2$ and $\sqrt3$ are badly
approximable. Markov's Theorem~\ref{ERG:thm:markov} is a detailed study of exactly these binary forms.

\begin{theorem}[Ratner 1991]\label{ERG:thm:ratner}
Let $G$ be a Lie group, $\Gamma\subset G$ a lattice and $U\subseteq G$ a subgroup generated by
one-parameter unipotent subgroups. Then the closure of every orbit $Ux\Gamma$ in $G/\Gamma$ is
\emph{homogeneous}: it is an orbit $Fx\Gamma$ of a closed subgroup $F\supseteq U$, carrying a finite
$F$-invariant measure. Every ergodic $U$-invariant probability measure is of this algebraic form.
\end{theorem}

This is the non-abelian descendant of Theorem~\ref{ERG:thm:kroneckerweyl}, where orbit closures of linear
flows were subtori. For diagonal flows nothing like it holds in general (bounded orbits abound, as the
badly approximable numbers show). But for higher rank diagonal actions there is a rigidity conjecture of
Margulis, of which the Littlewood conjecture in the Gallery is a consequence.

%======================================================================
\section{Gallery of theorems}
\label{ERG:sec:gallery}

\begin{itemize}[leftmargin=1.2em,itemsep=0.45em]
\item \textbf{Weyl's polynomial theorem} (Weyl 1916, \emph{Math.\ Ann.}~77). If a real polynomial $p$
  has an irrational coefficient other than the constant term, then $\{p(k)\}$ is equidistributed. This is
  the first step beyond rotations; the proof uses Weyl differencing.
\item \textbf{Bounded remainder intervals} (Hecke 1922, Ostrowski; Kesten 1966, \emph{Acta Arith.}~12,
  193--212). For irrational $\alpha$, the discrepancy $\#\{k<N:\{k\alpha\}\in I\}-N\abs I$ stays bounded
  if and only if $\abs I\in\Z+\alpha\Z$. The window of the Fibonacci chain is such an interval, which is
  why its tile counts deviate from $N\varphi^{-1}$, $N\varphi^{-2}$ by $O(1)$.
\item \textbf{Khinchin's theorem} (Khinchin 1924, \emph{Math.\ Ann.}~92). Let $\psi\colon\N\to[0,\infty)$ be
  non-increasing. For almost every $\alpha$, $\norm{q\alpha}_\Z<\psi(q)$ has infinitely many solutions if
  $\sum\psi(q)=\infty$, and only finitely many if $\sum\psi(q)<\infty$. In particular almost every number
  is \emph{not} badly approximable, yet the badly approximable numbers have full Hausdorff dimension
  (Jarn\'ik 1928).
\item \textbf{The Duffin--Schaeffer conjecture} (Duffin--Schaeffer 1941; proved by Koukoulopoulos and
  Maynard, \emph{Ann.\ of Math.}~192 (2020), 251--307). For \emph{any} $\psi\colon\N\to[0,\infty)$, if
  $\sum_q\psi(q)\phi(q)/q=\infty$ (with $\phi$ Euler's function), then almost every $\alpha$ has infinitely
  many reduced fractions $a/q$ with $\abs{\alpha-a/q}\le\psi(q)/q$. This removes monotonicity from
  Khinchin's theorem. The proof is a delicate study of overlaps of the approximation sets, organised as a
  graph-theoretic compression argument.
\item \textbf{Littlewood's conjecture} (c.~1930, open). For all real $\alpha,\beta$,
  $\liminf_{q\to\infty}q\norm{q\alpha}_\Z\norm{q\beta}_\Z=0$. Einsiedler, Katok and Lindenstrauss
  (\emph{Ann.\ of Math.}~164 (2006), 513--560) classified the positive entropy invariant measures for the
  diagonal group on $\SL_k(\R)/\SL_k(\Z)$, and deduced that the set of exceptions $(\alpha,\beta)$ has
  Hausdorff dimension $0$. By the Dani correspondence, Littlewood concerns orbits of the $2$-parameter
  diagonal group on $X_3$.
\item \textbf{Three gaps and lattices} (Marklof--Str\"ombergsson, \emph{Amer.\ Math.\ Monthly}~124 (2017),
  741--745). A proof of Theorem~\ref{ERG:thm:threegap} in which the gaps are the values of a function on the
  space of $2$-dimensional lattices, as in the remark at the end of Section~\ref{ERG:sec:threegap}.
\item \textbf{Gaps of $\sqrt n$ mod 1} (Elkies--McMullen, \emph{Duke Math.\ J.}~123 (2004), 95--139). The
  points $\{\sqrt n\}$, $n\le N$, also cut the circle (Lecture~7). Their normalized gap distribution
  exists, is \emph{not} the exponential law of random points, and is given by an explicit piecewise analytic
  formula. The proof uses Ratner's theorem on a space of affine lattices.
\item \textbf{Optimal dispersion} (Breneis--Hinrichs, arXiv 2019; in \emph{Discrepancy Theory}, de Gruyter
  2020, 117--132). The largest empty axis-parallel box among $n\ge2$ points in the torus $\T^2$ has area at least $2/n$,
  and for $n$ a Fibonacci number the Fibonacci lattice attains exactly $2/n$. This is a packing problem
  for boxes in which the golden lattice is optimal.
\item \textbf{Spherical Fibonacci points} (Aistleitner--Brauchart--Dick, \emph{Discrete Comput.\ Geom.}~48
  (2012), 990--1024). The spherical cap discrepancy is $O(N^{-1/2})$, with an explicit constant, for the spherical
  Fibonacci lattice, and the same order for spherical digital nets. Compare Beck's lower bound $cN^{-3/4}$
  (\emph{Mathematika}~31, 1984).
\item \textbf{Pure point diffraction} (Hof, \emph{Comm.\ Math.\ Phys.}~169 (1995), 25--43; Schlottmann,
  in \emph{Directions in Mathematical Quasicrystals}, CRM Monogr.~13, AMS 2000, 143--159). Regular model
  sets have pure point diffraction. This is the precise sense in which a model set is ``crystalline
  without being periodic''.
\item \textbf{Gauss circle problem for Penrose tilings} (Haynes--Lutsko, arXiv:2512.21444, 2025). The number of
  vertices of a rhombic Penrose tiling (unit edges) in a disc of radius $R$ is
  $\pi C_PR^2+O(R^{2/3}(\log R)^{2/3})$, uniformly over all such tilings, with $C_P\approx1.231$. The proof
  combines the cut and project description with Poisson summation in $\Q(\zeta_5)$, a direct bridge to
  Interlude~B. A recent preprint.
\item \textbf{Undecidability of tiling} (Wang 1961; Berger, \emph{Mem.\ Amer.\ Math.\ Soc.}~66 (1966);
  Robinson, \emph{Invent.\ Math.}~12 (1971)). It is undecidable whether a finite set of Wang tiles (unit
  squares with coloured edges, matched edge to edge, no rotations) tiles the plane. The proof needs, and
  produces, an \emph{aperiodic} tile set: one that tiles, but only non-periodically. Berger used $20426$
  tiles, and Robinson used $6$ tiles (allowing rotations and reflections). The minimum is $11$ Wang tiles, on $4$ colours (Jeandel--Rao,
  \emph{Advances in Combinatorics} 2021).
\item \textbf{Penrose tilings} (Penrose 1974, \emph{Bull.\ Inst.\ Math.\ Appl.}~10; de Bruijn 1981,
  \emph{Indag.\ Math.}~43). Two rhombi with matching rules tile only aperiodically. Every finite patch
  recurs with bounded gaps, and the vertex sets are model sets (Section~\ref{ERG:sec:cutproject}).
\item \textbf{The einstein problem} (Smith--Myers--Kaplan--Goodman-Strauss, ``An aperiodic monotile'',
  \emph{Combinatorial Theory}~4(1) (2024); ``A chiral aperiodic monotile'', \emph{ibid.}~4(2) (2024)). The
  ``hat'', a polygon made of $8$ kites of the deltoidal trihexagonal tiling, tiles the plane but never
  periodically. Its tilings use the hat and its mirror image. The ``spectre'', a modification of a related
  tile, tiles only aperiodically and needs no mirror images. Earlier, Socolar and Taylor (\emph{J.\ Combin.\ Theory Ser.~A}~118 (2011), 2207--2231) had
  given an aperiodic hexagonal monotile whose matching rules are not enforced by shape alone.
\item \textbf{Keller's conjecture} (Keller 1930). In every tiling of $\R^n$ by unit cubes, some two cubes
  share a full facet. It is true for $n\le6$ (Perron 1940). It is false for $n\ge10$ (Lagarias--Shor,
  \emph{Bull.\ Amer.\ Math.\ Soc.}~27 (1992), 279--283) and for $n\ge8$ (Mackey, \emph{Discrete Comput.\
  Geom.}~28 (2002), 275--279). It is true for $n=7$ (Brakensiek--Heule--Mackey--Narv\'aez, IJCAR 2020;
  \emph{J.\ Automated Reasoning}~66 (2022), 277--300, by a SAT solver). By Haj\'os' 1942 theorem the
  \emph{lattice} case, Minkowski's conjecture, is true in all dimensions. By a reduction of Szab\'o (1986)
  one may assume that the tiling is periodic modulo $2\Z^n$, so that it is a configuration of cubes in the
  flat torus $\R^n/2\Z^n$.
\item \textbf{Periodic tiling conjecture} (Grünbaum--Shephard; Lagarias--Wang 1996). If a single finite set
  $F\subset\Z^d$ tiles $\Z^d$ by translations, does it also tile periodically? Yes for $d=2$ (Bhattacharya,
  \emph{Amer.\ J.\ Math.}~142 (2020), 255--266), hence tiling by one tile is decidable there. No in general: Greenfeld and
  Tao (\emph{Ann.\ of Math.}~200 (2024), 301--363) found a counterexample in sufficiently high dimension,
  encoding a ``Sudoku'' by a single tiling equation. So translational tilings by one tile can be forced to
  be aperiodic.
\item \textbf{Convex pentagons} (Reinhardt 1918 to Mann--McLoud-Mann--Von Derau, arXiv 2015,
  \emph{Geom.\ Dedicata}~194 (2018), 141--167). Fifteen families of
  convex pentagons are known to tile the plane, the fifteenth found by computer search in 2015. Rao (arXiv
  2017, computer-assisted) announced a proof that the list is complete; Hales independently reproduced
  the main part of the computation, but to our knowledge the proof has not appeared in a refereed journal. As a consequence, no convex polygon
  tiles the plane only non-periodically.
\item \textbf{Oppenheim and Ratner} (Margulis, \emph{C.\ R.\ Acad.\ Sci.\ Paris}~304 (1987); Ratner,
  \emph{Ann.\ of Math.}~134 (1991), 545--607, and \emph{Duke Math.\ J.}~63 (1991), 235--280). See
  Theorems~\ref{ERG:thm:margulis} and~\ref{ERG:thm:ratner}. They are the reason that counting lattice
  points in thin regions (Interlude~B) can be done by dynamics.
\item \textbf{Quasicrystals and icosians} (Elser--Sloane, \emph{J.\ Phys.\ A}~20 (1987), 6161; Moody--Patera,
  \emph{J.\ Phys.\ A}~26 (1993), 2829). These are the $E_8\to H_4$ projection and its algebraic form, via the
  icosian ring. Shcherbak (\emph{Russian Math.\ Surveys}~43 (1988), 149) found the foldings
  $A_4\to H_2$, $D_6\to H_3$, $E_8\to H_4$ in singularity theory. Affine versions: Patera--Twarock,
  \emph{J.\ Phys.\ A}~35 (2002), 1551--1574; Dechant--Boehm--Twarock, \emph{J.\ Phys.\ A}~45 (2012), 285202
  and \emph{J.\ Math.\ Phys.}~54 (2013), 093508. Viruses: Twarock, \emph{J.\ Theor.\ Biol.}~226 (2004),
  477--482; Keef--Twarock, \emph{J.\ Math.\ Biol.}~59 (2009), 287--313.
\item \textbf{Quasicrystals in the lab} (Shechtman--Blech--Gratias--Cahn, \emph{Phys.\ Rev.\ Lett.}~53 (1984),
  1951; Nobel Prize in Chemistry 2011). Sharp icosahedral diffraction in Al--Mn: the experimental
  counterpart of pure point diffraction of model sets.
\end{itemize}

%======================================================================
\section*{Exercises}
\addcontentsline{toc}{section}{Exercises}

\begin{exercise}
Let $\gcd(p,q)=1$. Show that the closed geodesic $\ell_{p/q}$ on $\R^2/\Z^2$ has length $\sqrt{p^2+q^2}$,
and that its complement is a single open annulus of width $1/\sqrt{p^2+q^2}$. (Its lifts to $\R^2$ are the
lines $qy-px\in\Z$.) So rational lines are spread evenly, but they leave a gap that shrinks only like
$1/\text{length}$.
\end{exercise}

\begin{exercise}
Prove Theorem~\ref{ERG:thm:kroneckerweyl} for $\Lat=\Z^n$ using Weyl's criterion on $\T^n$ (characters
$e(\ip{m}{x})$, $m\in\Z^n$). Then deduce the general case by a linear change of variables, and identify
the dual lattice.
\end{exercise}

\begin{exercise}
For $\alpha=\varphi^{-1}$ and $N=F_n$, list the points $\{k\alpha\}$ for $n=5$ ($N=5$) and $n=6$
($N=8$), and verify Example~\ref{ERG:ex:fibgaps} directly. What happens for $N=F_n+1$?
\end{exercise}

\begin{exercise}\label{ERG:ex:oppenheim2}
Show that $Q(x,y)=(x-\sqrt2y)(x-\sqrt3y)$ satisfies $\inf_{(x,y)\in\Z^2\setminus0}\abs{Q(x,y)}>0$.
\emph{Hint:} $\sqrt2=[1;2,2,\dots]$ and $\sqrt3=[1;1,2,1,2,\dots]$ have bounded partial quotients. If
$\abs{x-\sqrt2y}<\frac{\sqrt3-\sqrt2}{2}\abs y$, then the other factor is at least
$\frac{\sqrt3-\sqrt2}{2}\abs y$.
\end{exercise}

\begin{exercise}
Show that the Fibonacci word $f$ is \emph{balanced}: any two factors of the same length contain numbers of
$b$'s differing by at most one. Deduce again that $bb$ does not occur, and show that $aaa$ does not occur
either. Interpret both facts for the jammed channel of Lecture~11.
\end{exercise}

\begin{exercise}
Verify the edge rule of Section~\ref{ERG:sec:fold} for $H_3$: the resulting graph on the
six nodes $\alpha_i,\beta_i$ is $D_6$. Then compute the $2\times2$ Gram matrix
of $a_1,a_2$ in the case $A_4\to H_2$ and check that the angle is $4\pi/5$.
\end{exercise}

\begin{exercise}
Using the explicit $A_4$ roots $e_i-e_j\in\R^5$, find the plane on which the cyclic shift
$e_1\to e_2\to\dots\to e_5\to e_1$ acts as rotation by $2\pi/5$, project the $20$ roots to it, and recover
Figure~\ref{ERG:fig:a4}. (The cyclic shift is a Coxeter element of $W(A_4)=\Sigma_5$.)
\end{exercise}

\begin{exercise}
Using only Siegel's mean value theorem, show that for every $\varepsilon>0$ there is a lattice in $\R^n$
with packing density at least $(2-\varepsilon)2^{-n}$. Where exactly is the symmetry $v\leftrightarrow-v$
used?
\end{exercise}

\begin{exercise}\label{ERG:ex:mh}
Prove the first half of the Morse--Hedlund theorem. If $p(n)\le n$ for some $n$, then $p(m)=p(m+1)$ for
some $m\le n$. Deduce that every factor of length $m$ has a unique extension to the right, so the word is
eventually periodic.
\end{exercise}

\section*{Further reading}
\addcontentsline{toc}{section}{Further reading}
\begin{itemize}[nosep]
\item M.~Einsiedler and T.~Ward, \emph{Ergodic Theory with a view towards Number Theory}, Springer GTM 259,
  2011. Rotations, Weyl, the space of lattices, and the Dani correspondence.
\item J.~W.~S.~Cassels, \emph{An Introduction to Diophantine Approximation}, Cambridge, 1957; and
  M.~Aigner, \emph{Markov's Theorem and 100 Years of the Uniqueness Conjecture}, Springer, 2013.
\item L.~Kuipers and H.~Niederreiter, \emph{Uniform Distribution of Sequences}, Wiley, 1974. Weyl's
  criterion, discrepancy, $\{k\alpha\}$.
\item M.~Lothaire, \emph{Algebraic Combinatorics on Words}, Cambridge, 2002, Ch.~2 (Sturmian words).
\item M.~Baake and U.~Grimm, \emph{Aperiodic Order}, Vol.~1, Cambridge, 2013. Substitutions, cut and
  project, Penrose tilings, diffraction.
\item M.~Senechal, \emph{Quasicrystals and Geometry}, Cambridge, 1995; B.~Gr\"unbaum and G.~C.~Shephard,
  \emph{Tilings and Patterns}, Freeman, 1987.
\item J.~E.~Humphreys, \emph{Reflection Groups and Coxeter Groups}, Cambridge, 1990, for $H_3$, $H_4$ and
  Coxeter elements. Conway--Sloane, Ch.~8, for the icosians.
\item D.~W.~Morris, \emph{Ratner's Theorems on Unipotent Flows}, Chicago Lectures in Math., 2005.
\end{itemize}
